% Options for packages loaded elsewhere
% Options for packages loaded elsewhere
\PassOptionsToPackage{unicode}{hyperref}
\PassOptionsToPackage{hyphens}{url}
\PassOptionsToPackage{dvipsnames,svgnames,x11names}{xcolor}
%
\documentclass[
  german,
  10pt,
  twoside]{article}
\usepackage{xcolor}
\usepackage[top=2.5cm,bottom=2.5cm,left=2cm,right=2cm]{geometry}
\usepackage{amsmath,amssymb}
\setcounter{secnumdepth}{5}
\usepackage{iftex}
\ifPDFTeX
  \usepackage[T1]{fontenc}
  \usepackage[utf8]{inputenc}
  \usepackage{textcomp} % provide euro and other symbols
\else % if luatex or xetex
  \usepackage{unicode-math} % this also loads fontspec
  \defaultfontfeatures{Scale=MatchLowercase}
  \defaultfontfeatures[\rmfamily]{Ligatures=TeX,Scale=1}
\fi
\usepackage{lmodern}
\ifPDFTeX\else
  % xetex/luatex font selection
\fi
% Use upquote if available, for straight quotes in verbatim environments
\IfFileExists{upquote.sty}{\usepackage{upquote}}{}
\IfFileExists{microtype.sty}{% use microtype if available
  \usepackage[]{microtype}
  \UseMicrotypeSet[protrusion]{basicmath} % disable protrusion for tt fonts
}{}
\usepackage{setspace}
\makeatletter
\@ifundefined{KOMAClassName}{% if non-KOMA class
  \IfFileExists{parskip.sty}{%
    \usepackage{parskip}
  }{% else
    \setlength{\parindent}{0pt}
    \setlength{\parskip}{6pt plus 2pt minus 1pt}}
}{% if KOMA class
  \KOMAoptions{parskip=half}}
\makeatother
% Make \paragraph and \subparagraph free-standing
\makeatletter
\ifx\paragraph\undefined\else
  \let\oldparagraph\paragraph
  \renewcommand{\paragraph}{
    \@ifstar
      \xxxParagraphStar
      \xxxParagraphNoStar
  }
  \newcommand{\xxxParagraphStar}[1]{\oldparagraph*{#1}\mbox{}}
  \newcommand{\xxxParagraphNoStar}[1]{\oldparagraph{#1}\mbox{}}
\fi
\ifx\subparagraph\undefined\else
  \let\oldsubparagraph\subparagraph
  \renewcommand{\subparagraph}{
    \@ifstar
      \xxxSubParagraphStar
      \xxxSubParagraphNoStar
  }
  \newcommand{\xxxSubParagraphStar}[1]{\oldsubparagraph*{#1}\mbox{}}
  \newcommand{\xxxSubParagraphNoStar}[1]{\oldsubparagraph{#1}\mbox{}}
\fi
\makeatother


\usepackage{graphicx}
\makeatletter
\newsavebox\pandoc@box
\newcommand*\pandocbounded[1]{% scales image to fit in text height/width
  \sbox\pandoc@box{#1}%
  \Gscale@div\@tempa{\textheight}{\dimexpr\ht\pandoc@box+\dp\pandoc@box\relax}%
  \Gscale@div\@tempb{\linewidth}{\wd\pandoc@box}%
  \ifdim\@tempb\p@<\@tempa\p@\let\@tempa\@tempb\fi% select the smaller of both
  \ifdim\@tempa\p@<\p@\scalebox{\@tempa}{\usebox\pandoc@box}%
  \else\usebox{\pandoc@box}%
  \fi%
}
% Set default figure placement to htbp
\def\fps@figure{htbp}
\makeatother



\ifLuaTeX
\usepackage[bidi=basic]{babel}
\else
\usepackage[bidi=default]{babel}
\fi
% get rid of language-specific shorthands (see #6817):
\let\LanguageShortHands\languageshorthands
\def\languageshorthands#1{}
\ifLuaTeX
  \usepackage[german]{selnolig} % disable illegal ligatures
\fi


\setlength{\emergencystretch}{3em} % prevent overfull lines

\providecommand{\tightlist}{%
  \setlength{\itemsep}{0pt}\setlength{\parskip}{0pt}}



 
\usepackage[style=authoryear,backend=biber,style=authoryear-comp,uniquename=false,uniquelist=false]{biblatex}
\addbibresource{references.bib}


\usepackage{booktabs}
\usepackage{longtable}
\usepackage{array}
\usepackage{multirow}
\usepackage{wrapfig}
\usepackage{float}
\usepackage{colortbl}
\usepackage{pdflscape}
\usepackage{tabu}
\usepackage{threeparttable}
\usepackage{threeparttablex}
\usepackage[normalem]{ulem}
\usepackage{makecell}
\usepackage{xcolor}
\usepackage{csquotes}
\usepackage{booktabs}
\usepackage{array}
\usepackage{multirow}
\usepackage{wrapfig}
\usepackage{float}
\usepackage{colortbl}
\usepackage{pdflscape}
\usepackage{threeparttable}
\usepackage{threeparttablex}
\usepackage[normalem]{ulem}
\usepackage{makecell}
\usepackage{xcolor}
\usepackage{multicol}
\usepackage{fancyhdr}
\usepackage{titlesec}
\usepackage{etoolbox}
\usepackage{tabularx}
\usepackage{afterpage}
\usepackage{tikz}
\usepackage{longtable}



% Suppress default maketitle
\renewcommand{\maketitle}{}

% Define WISTA colors (Havelock Blue as main)
\definecolor{wistaBlue}{RGB}{65,125,180}
\definecolor{wistaLightBlue}{RGB}{120,170,215}
\definecolor{wistaDarkGray}{RGB}{64,64,64}
\definecolor{wistaLightGray}{RGB}{128,128,128}

% Section formatting (##) - number 20% bigger, lines closer, text closer to bottom line
% e.g., "2" with lines around "Einleitung"
\titleformat{\section}[display]
  {\normalfont\large\bfseries}
  {\color{wistaBlue}\Large\bfseries\thesection\\\noindent\rule{\columnwidth}{1.5pt}}
  {0pt}
  {\vspace{0pt}\color{wistaDarkGray}\large\bfseries}
  [\vspace{0pt}\noindent{\color{wistaBlue}\rule{\columnwidth}{1.5pt}}\vspace{0.05cm}]
\titlespacing*{\section}{0pt}{1.5ex plus 0.3ex minus .1ex}{0.2ex plus .1ex}

% Subsection formatting (###) - grey text with blue line under
% e.g., "2.1 Handlungsbedarf"
\renewcommand{\thesubsection}{\thesection.\arabic{subsection}}
\titleformat{\subsection}
  {\color{wistaDarkGray}\normalsize\bfseries}
  {\color{wistaDarkGray}\normalsize\bfseries\thesubsection}
  {0.5em}
  {\color{wistaDarkGray}\normalsize\bfseries}
  [\vspace{0pt}\noindent{\color{wistaBlue}\rule{\columnwidth}{0.8pt}}\vspace{0.05cm}]
\titlespacing*{\subsection}{0pt}{1ex plus 0.2ex minus .1ex}{0.2ex plus .1ex}
  
% Subsubsection formatting (####) - grey text, no underline
% e.g., "2.1.1 Steigende Datenmenge"
\renewcommand{\thesubsubsection}{\thesubsection.\arabic{subsubsection}}
\titleformat{\subsubsection}
  {\color{wistaDarkGray}\normalsize\bfseries}
  {\color{wistaDarkGray}\normalsize\bfseries\thesubsubsection}
  {0.5em}
  {\color{wistaDarkGray}\normalsize\bfseries}

% Paragraph formatting (4th level) - numbered as subsection.subsubsection.paragraph
\renewcommand{\theparagraph}{\thesubsubsection.\arabic{paragraph}}
\titleformat{\paragraph}[runin]
  {\color{wistaDarkGray}\normalsize\bfseries}
  {\color{wistaBlue}\normalsize\bfseries\theparagraph}
  {0.5em}
  {\color{wistaDarkGray}\normalsize\bfseries}
  [.]
\titlespacing*{\paragraph}{0pt}{1ex plus 0.5ex minus 0.2ex}{0.5em}

% Header and footer setup with blue line under header
\pagestyle{fancy}
\fancyhf{}
\fancyhead[L]{\small\textit{Oliver Hauke, Annette Kristiansen}}
\fancyhead[R]{\small\textit{Effiziente Analyse großer Forschungsdaten}}
\fancyfoot[L]{\small\thepage}
\fancyfoot[R]{\small Statistisches Bundesamt | WISTA | {\color{wistaBlue}6} | 2025}
\renewcommand{\headrulewidth}{0pt}
\renewcommand{\footrulewidth}{0pt}
% Add blue line under header
\renewcommand{\headrule}{\color{wistaBlue}\hrule width\headwidth height 1pt \vskip 2pt}
\renewcommand{\headrulewidth}{1pt}

% Make all figures span both columns (full width)
\let\origfigure\figure
\let\endorigfigure\endfigure
\renewenvironment{figure}[1][htbp]{%
  \origfigure*[#1]%
}{%
  \endorigfigure*%
}

% Allow manual control of table* environment for full-width tables
% Tables by default stay in single column, but can use table* manually



% Make bibliography full-width (one column) - for biblatex
\defbibheading{bibliography}[\bibname]{%
  \onecolumn
  \section*{#1}%
  \markboth{#1}{#1}%
}



% Footnotes stay within column
\usepackage[stable]{footmisc}

% Custom footnote formatting: "|^1" in dark WISTA gray
\renewcommand{\thefootnote}{{\color{wistaDarkGray}|$^{\arabic{footnote}}$}}

% Pifont for dingbat symbols
\usepackage{pifont}

% Better table handling for two-column mode
\usepackage{array,tabularx,calc}

% Custom WISTA table format
\usepackage{caption}
\usepackage{xcolor}
\usepackage{colortbl}
\usepackage{array}

% Define custom caption format for WISTA tables
\DeclareCaptionFormat{wista}{%
  {\color{wistaLightBlue}\bfseries #1#2}\par\vspace{0.2em}%
  {\color{wistaDarkGray}\normalfont #3}\par\vspace{0.3em}%
}

% Set up table captions with custom WISTA format
\captionsetup[table]{
  position=top,
  format=wista,
  justification=justified,
  singlelinecheck=false,
  labelsep=none,
  skip=0.3em
}

% Better table spacing and white inner lines
\setlength{\tabcolsep}{6pt}
\renewcommand{\arraystretch}{1.15}

% Set white lines for tables by default, blue for first column
\AtBeginDocument{
  \arrayrulecolor{white}
  \setlength{\arrayrulewidth}{0.3pt}
}

% Define commands for colored table borders
\newcommand{\bluevrule}{\color{wistaBlue}\vrule}
\newcommand{\whitevrule}{\color{white}\vrule}

% Define WISTA table colors
\definecolor{wistaTableBg}{RGB}{240,248,255}     % Very light blue background
\definecolor{wistaHeaderBg}{RGB}{220,230,240}    % Light blue header background
\definecolor{wistaFirstCol}{RGB}{65,125,180}     % Blue for first column lines

% Custom WISTA table command
\newcommand{\wistaTableStart}{%
  \rowcolors{2}{wistaTableBg}{wistaTableBg}%
  \arrayrulecolor{wistaFirstCol}%
}

% Custom commands for table styling
\newcommand{\wistaHeaderRow}[1]{%
  \rowcolor{wistaHeaderBg}%
  \arrayrulecolor{wistaDarkGray}%
  #1 \\
  \arrayrulecolor{white}%
}

\newcommand{\wistaFirstColRule}{\arrayrulecolor{wistaFirstCol}}
\newcommand{\wistaWhiteRule}{\arrayrulecolor{white}}

% Custom bullet points with > character - 35% width, 95% height, black extra bold, half line spacing
\usepackage{enumitem}
\setlist[itemize]{label={\scalebox{0.35}[0.95]{\fontseries{bx}\selectfont >}}, leftmargin=1.5em, itemsep=5pt, topsep=3pt, parsep=0pt}

% Redefine ref to include pifont 247 (arrow) for figure/table/section references
\let\oldref\ref
\renewcommand{\ref}[1]{\oldref{#1}\raisebox{-0.2ex}{\scriptsize\color{wistaBlue}\ding{247}}}
\makeatletter
\@ifpackageloaded{caption}{}{\usepackage{caption}}
\AtBeginDocument{%
\ifdefined\contentsname
  \renewcommand*\contentsname{Inhaltsverzeichnis}
\else
  \newcommand\contentsname{Inhaltsverzeichnis}
\fi
\ifdefined\listfigurename
  \renewcommand*\listfigurename{Abbildungsverzeichnis}
\else
  \newcommand\listfigurename{Abbildungsverzeichnis}
\fi
\ifdefined\listtablename
  \renewcommand*\listtablename{Tabellenverzeichnis}
\else
  \newcommand\listtablename{Tabellenverzeichnis}
\fi
\ifdefined\figurename
  \renewcommand*\figurename{Abbildung}
\else
  \newcommand\figurename{Abbildung}
\fi
\ifdefined\tablename
  \renewcommand*\tablename{Tabelle}
\else
  \newcommand\tablename{Tabelle}
\fi
}
\@ifpackageloaded{float}{}{\usepackage{float}}
\floatstyle{ruled}
\@ifundefined{c@chapter}{\newfloat{codelisting}{h}{lop}}{\newfloat{codelisting}{h}{lop}[chapter]}
\floatname{codelisting}{Listing}
\newcommand*\listoflistings{\listof{codelisting}{Listingverzeichnis}}
\makeatother
\makeatletter
\makeatother
\makeatletter
\@ifpackageloaded{caption}{}{\usepackage{caption}}
\@ifpackageloaded{subcaption}{}{\usepackage{subcaption}}
\makeatother
\usepackage{bookmark}
\IfFileExists{xurl.sty}{\usepackage{xurl}}{} % add URL line breaks if available
\urlstyle{same}
\hypersetup{
  pdftitle={EFFIZIENTE AUSWERTUNG DES BUSINESS-TAX-PANELS MITHILFE VON R},
  pdfauthor={Oliver Hauke; Annette Kristiansen},
  pdflang={de},
  colorlinks=true,
  linkcolor={blue},
  filecolor={Maroon},
  citecolor={Blue},
  urlcolor={Blue},
  pdfcreator={LaTeX via pandoc}}


\title{EFFIZIENTE AUSWERTUNG DES BUSINESS-TAX-PANELS MITHILFE VON R}
\author{Oliver Hauke \and Annette Kristiansen}
\date{2025-12-05}
\begin{document}
\maketitle

% Custom WISTA-style title page
\begin{titlepage}
\thispagestyle{empty}
\pagestyle{empty}
\onecolumn

% Horizontal blue line at top (header line)
\noindent{\color{wistaBlue}\rule{\textwidth}{1pt}}

\vspace{-0.5cm}

% Vertical blue line on the left (6.3cm from left edge, starts below header line, stops above footer)
\begin{tikzpicture}[remember picture, overlay]
  \draw[wistaBlue, line width=2pt] 
    ([xshift=6.3cm, yshift=-4cm]current page.north west) -- 
    ([xshift=6.3cm, yshift=3.5cm]current page.south west);
\end{tikzpicture}

\vspace*{0.5cm}

% Left side - short author bios (right-aligned to the vertical line)
\begin{minipage}[t]{3.9cm}
\selectlanguage{ngerman}%
\raggedleft
\hyphenpenalty=0
\tolerance=1000
\fontsize{7}{8.5}\selectfont
\vspace{0pt}%
{\fontsize{8}{9.5}\selectfont\color{wistaBlue}\textbf{Oliver Hauke}}\\[0.15cm]
ist Mathematiker und IT-Referent im Kompetenzzentrum „Analyse und Auswertung" des Statistischen Bundesamtes. Er verantwortet die Weiterentwicklung der technischen Infrastruktur des Forschungsdatenzentrums und berät das Netzwerk empirische Steuerforschung in der effizienten Nutzung der Systeme.\\[0.5cm]

{\fontsize{8}{9.5}\selectfont\color{wistaBlue}\textbf{Annette Kristiansen}}\\[0.15cm]
ist Ökonomin und Referentin im Referat „Unternehmenssteuern" des Statistischen Bundesamtes. Sie beschäftigt sich mit der Zusammenführung der Unternehmenssteuerstatistiken und betreut das Netzwerk empirische Steuerforschung. Für ihre externe Promotion an der Universität Bremen untersucht sie die technologische Vielfalt multinationaler Unternehmen.

\end{minipage}%
\hspace{0.4cm}% Gap to vertical line (left side)
\hspace{0.4cm}% Gap from vertical line (right side)
% Right side - title, authors, keywords, abstracts
\begin{minipage}[t]{11.8cm}
\raggedright
\vspace{1.2cm}%

% Title (uppercase via command, not bold)
{\LARGE\color{wistaDarkGray}\begin{spacing}{1.3}\MakeUppercase{Effiziente Auswertung des Business-Tax-Panels mithilfe von R}\end{spacing}}

\vspace{-0.2cm}

% Blue horizontal line (underline for title)
{\color{wistaBlue}\rule{\linewidth}{0.5pt}}

\vspace{0.5cm}

% Authors inline, dark gray
{\large\color{wistaDarkGray} Oliver Hauke, Annette Kristiansen}

\vspace{0.05cm}

% Blue line (underline for authors)
{\color{wistaBlue}\rule{\linewidth}{0.5pt}}

\vspace{0.5cm}

% Keywords (pifont 247 arrow then bold label and blue text)
\hangindent=1.2em\hangafter=1
{\color{wistaBlue}\ding{247}\ {\bfseries Schlüsselwörter:} Effiziente Analysen, R-Programmierung, Forschungsdaten, Unternehmenssteuerstatistik, Business-Tax-Panel, Netzwerk empirische Steuerforschung}

\vspace{0.3cm}

% Zusammenfassung (uppercase via command)
{\color{wistaDarkGray}\bfseries\MakeUppercase{Zusammenfassung}}\\[0.15cm]
\small
Das Forschungsdatenzentrum des Statistischen Bundesamtes baut seine technische Infrastruktur gezielt aus, um der Wissenschaft erweiterte Analysemöglichkeiten amtlicher Daten bereitzustellen. Seit November 2025 steht Forschenden exklusiver Zugang zu diesen erweiterten Systemressourcen in R und Stata zur Verfügung. R-basierte Tools -- insbesondere Apache Arrow, Parquet und DuckDB -- ermöglichen es, gängige Analysen großer Forschungsdatensätze in Echtzeit auszuführen. Am Beispiel des Business-Tax-Panel (2013 bis 2019) werden die erzielbaren Effizienzgewinne bei der Auswertung umfangreicher steuerstatistischer Daten im Rahmen des Netzwerks empirische Steuerforschung demonstriert.

\vspace{0.5cm}

% English keywords and abstract (italicized)
\hangindent=1.2em\hangafter=1
{\itshape\color{wistaBlue}\ding{247}\ {\bfseries Keywords:} Efficient analysis -- R programming -- research data -- corporate tax statistics -- Business Tax Panel -- empirical tax research network}

\vspace{0.3cm}

{\itshape\color{wistaDarkGray}\bfseries\MakeUppercase{Abstract}}\\[0.15cm]
\small\itshape
The Research Data Centre at the Federal Statistical Office is strategically expanding its technical infrastructure to provide researchers with enhanced analytical capabilities for official data. Since November 2025, researchers have had exclusive access to these extended system resources in R and Stata. R-based tools -- particularly Apache Arrow, Parquet, and DuckDB -- enable real-time analysis of large research datasets. Using the Business Tax Panel (2013--2019), this paper demonstrates the achievable efficiency gains in analysing large-scale tax statistics within the empirical tax research network.

\end{minipage}

\vfill

\vspace{0.5cm}
% Footer matching document style
\noindent\makebox[\textwidth]{%
  \small 1%
  \hfill%
  \small Statistisches Bundesamt | WISTA | {\color{wistaBlue}6} | 2025%
}

\end{titlepage}

% Stay in one-column mode for TOC (will be switched later)
\pagestyle{fancy}

\renewcommand*\contentsname{Inhaltsverzeichnis}
{
\hypersetup{linkcolor=}
\setcounter{tocdepth}{3}
\tableofcontents
}

\setstretch{1}
% Table of Contents will be inserted here by Quarto
\clearpage
% Switch to two-column mode for main document
\twocolumn

\section{Einleitung}\label{sec-einleitung}

\ldots{}

Spannungsfeld, Weiterentwicklung und Technologie-Empfehlung\ldots{}

Forschungsdatenzentrum des Statistischen Bundesamts, hier kurz FDZ
Netzwerk empririscher Steuerforschung, hier kurz NeSt

\subsection{Handlungsbedarf}\label{handlungsbedarf}

\subsubsection{Steigende Datenmenge in der amtlichen
Statistik}\label{steigende-datenmenge-in-der-amtlichen-statistik}

\begin{itemize}
\item
  Qualitätsbericht. Wachsende Datenmenge in der amtl. Statistik, Verweis
  auf \autocite{destatis2023}
\item
  Datenverarbeitung im StBA \ldots{}
\item
  neues NeSt Produkt BTP zeigt Umfang, hoch interessant für Wissenschaft
\item
  Datenverarbeitungsaufgaben des StBA erklären
\item
  Nicht mehr durch vorhandene Systemkapazitäten (insb. Arbeitsspeicher)
  gedeckt
\item
  Herausforderungen (größere Datenmengen, komplexere Analysen,
  schnellere Antworten notwendig)

  \begin{itemize}
  \tightlist
  \item
    hohe Anforderungen der Öffentlichkeit und besonders Wissenschaft.
    Verweis auf Gutachten des Stat. Beirats und Gründung ``Netzwerk emp.
    Steuerforschung''

    \begin{itemize}
    \tightlist
    \item
      Schnelle Bereitstellung von am besten vollständigen, wenig
      Anonymisierten Datensätzen an Wissenschaft via FDZ. Erste Analysen
      seitens StBA und hohe Expertise erwartet (Exploration der Daten
      nötig)
    \item
      Bereitstellung von Zusatzdokumentation durch StBA. Metadatenreport
    \item
      schnelle Bereitstellung,
    \item
      zeitnahe Nutzung in FDZ
    \item
      keine Wartezeiten in FDZ
    \item
      sofortig Antworten auf Analysefragen
    \item
      Erhöhte Komplexität Analysen mit größeren, Komplexeren Datensätzen
    \end{itemize}
  \item
    Rückstände/Altlasten: weiterentwickelnde Infrastruktur, nicht state
    of the art (ausgelegt auf konsistenz und zuverlässigkeit, nicht
    schnelligkeit), gesetzlicher Auftrag

    \begin{itemize}
    \tightlist
    \item
      hauptsächlich SAS für große Datenmenge verwendet, funktioniert,
      aber langsam
    \item
      FDZ Bund integriert in Standardarbeitsplätze des StBA
    \item
      Bereitstellungsaufwände und Abhängigkeiten moderner Software
    \end{itemize}
  \end{itemize}
\end{itemize}

Annette: -\textgreater{} Impuls aus der Wissenschaft, dann Fachbereich -
Fachbereich/Verbund im Grunde Hauptnutzende da Tools und Daten zur
Verfügung stehen Technische Möglichkeit entwickelt sich weiter usw..

Der Sachverständigenrat zur Begutachtung der gesamtwirtschaftlichen
Entwicklung (2023) hat in seinem Gutachten 2023/24 auf eine technisch
veraltete Infrastruktur und auf die fehlende Nutzerfreundlichkeit
hingewiesen. Seitdem hat sich die die Infrastruktur in dem
Forschungsdatenzentrum des Bundes verbessert. Seit Januar 2025 besteht
ein Remote Access System als neuer Zugangsweg für die Wissenschaft. Im
Bereich der Steuerstatistiken wird derzeit geprüft, inwiefern auch
Einzeldaten aus diesem Bereich über Remote Access genutzt werden können.
Als Übergang und Erleichterung für bestehende Forschungsprojekte wird in
diesem Kapitel 2

Die amtliche Statistik steht vor der Herausforderung, bei stetig
wachsenden Datenmengen und sich ändernden Nutzeranforderungen qualitativ
hochwertige und zeitnahe statistische Informationen bereitzustellen
\autocite{destatis2023}. Dabei spielt die Effizienz der verwendeten
statistischen Verfahren eine zentrale Rolle für die Leistungsfähigkeit
des gesamten statistischen Systems.

In den vergangenen Jahren haben sich die Rahmenbedingungen für die
Produktion amtlicher Statistiken erheblich verändert. Die
Digitalisierung hat nicht nur zu einer Vervielfachung der verfügbaren
Datenquellen geführt, sondern auch neue Möglichkeiten für die
statistische Analyse eröffnet. Gleichzeitig sind die Erwartungen der
Nutzer an Aktualität und Detailliertheit der Statistiken gestiegen.

Diese Entwicklungen erfordern eine kontinuierliche Evaluation und
Optimierung der eingesetzten Verfahren. Ziel der vorliegenden Studie ist
es daher, die Effizienz verschiedener statistischer Verfahren
systematisch zu bewerten und Handlungsempfehlungen für die Praxis
abzuleiten.

\subsubsection{Steigende Anforderungen der Öffentlichkeit insb.
Wissenschaft}\label{steigende-anforderungen-der-uxf6ffentlichkeit-insb.-wissenschaft}

FDZ, NeSt

\begin{itemize}
\tightlist
\item
  Schnelle Reaktionen erforderlich
\item
  Komplexere (Forschungsfragen) zu beantworten
\item
  FDZ erklären
\item
  Gutachten BMF ansprechen
\item
  NeSt erklären, Bedarfe durch Produkt BTP beantwortet.. aber wie kann
  es in angme
\end{itemize}

The Need for Speed

As we start working with larger and larger datasets, the basic tools of
the tidyverse start to look a little slow. In the last few years several
packages more suited to large datasets have emerged. Some of these are
column, rather than row, oriented. Some use parallel processing. Some
are vector optimized. Speedy databases that have made their way into the
R ecosystem are data.table, arrow, polars and duckdb. All of these are
available for Python as well. Each of these carries with it its own
interface and learning curve. duckdb, for example is a dialect of SQL,
an entirely different language so our dplyr code above has to look like
this in SQL: S.
https://www.r-bloggers.com/2024/04/the-truth-about-tidy-wrappers-3/

\subsection{Vorschreitende Innovation der
Datenverarbeitungstechnologien}\label{vorschreitende-innovation-der-datenverarbeitungstechnologien}

\subsubsection{Ausgangslage}\label{ausgangslage}

SAS für amtl. Statistik und SAS/Stata für FDZ

\begin{itemize}
\tightlist
\item
  SAS, zsl. Stata in den FDZ. Session zu SAS Ablösung
\item
  Weiterentwicklung der Infrastruktur: neue Server für R im StBA und für
  FDZ eigene neue Server für R und Stata
\item
  hohe Lizenzkosten, Trend zu Abonnements und Cloud nachteilig,
  technologielock und Privacy.
\item
  Kommen bei unseren Anwendungen Kapazitär an ihre Grenzen und
  Performanz in anderen Technologien höher.
\end{itemize}

Anforderungen an Modernisierung

\begin{itemize}
\tightlist
\item
  hohe Performanz, schnelle Ergebnisse, idealenfalls in Echtzeit (d.h.
  wenigen Sekunden)
\item
  moderne Interfaces, möglichst leichte Nutzung, wenig Schulungsbedarf,
  Allgemeinwissen von Statistiker u.Ä. Fachrichtungen (Werkzeug, nicht
  Gegenstand der Untersuchung)
\item
  leicht einrichtbar, wenig Abhängigkeiten zu vielen Tools, die sich
  ständig weiterentwicklen, sodass hohe Wartungsaufwände entstehen.
\end{itemize}

\subsubsection{Vormarsch von Open
Source}\label{vormarsch-von-open-source}

\begin{itemize}
\tightlist
\item
  in der internationalen amtlichen Statistik seit Jahren auf dem
  Vormarsch. Langsam steigt auch Verbreitung von R steigt auch in den
  dt. Statistikämtern.
\end{itemize}

\subsubsection{Fokus auf R}\label{fokus-auf-r}

R verbreitung steigt, bei uns in Kinderschuhen. Wie fängt man denn am
besten an? Was ist von anfang an der richtige Weg daten in R zu
verarbeiten?

R ist weltweit bekannt und verbreitet als die Lingua Franca für
Statistik, Methodologie und Data Science. Dies macht R zu der
natürlichen Wahl für die amtlichen Statistik mit viele attraktive
Vorteile.

\begin{itemize}
\item
  aktive globale Community sorgt für schnelle Innovation und
  Verbesserung
\item
  CRAN bietet offiziell aktuell Zusatzfunktionalitäten in Form von X
  R-Paketen. Darunter sind viele klassische Pakete für Datenverarbeitung
  und -anayse.
\item
  R

  \begin{itemize}
  \tightlist
  \item
    The reasons why the official statistics community is rapidly
    embracing R are clear. S. erfolg / Wachstum der Conference uRos
    https://r-project.ro/conference2025.html\#Presentations
  \item
    R lingua franca for statisticians, methodologists and data
    scientists worldwide -\textgreater{} nativ geeignet für amtl.
    Statistik
  \item
    R Open Source bietet vorteile

    \begin{itemize}
    \tightlist
    \item
      aktive globale Community
    \item
      support from the industry
    \item
      it combines a vast amount of functionalities for data preparation,
      methodology, visualisation and application building
    \item
      free = ideal environment for sharing, collaboration and
      industrialising official statistics
    \item
      Packages CRAN, github\ldots{}
    \end{itemize}
  \end{itemize}
\end{itemize}

Arrow bspw. hatte in 2024 und 2025 jeweils 4 Major Releases mit diversen
neuen Funktionalitäten und Performanz-Verbesserung

Neue Entwicklung zu R in der dt. OSS, bei Integration in
Statistikproduktion und Datenbereitstellung direkt mit den richtigen
Features und Methoden beginnen. Wie die folgenden\ldots{}

Open Source -\textgreater{} breite Auswahl an Parketen zum Einlesen,
Transformieren und Auswerten von Daten.

\begin{itemize}
\tightlist
\item
  Unterscheidung klassischer/r-nativer Methoden (base-r, readr / vroom
  und data.table), effiziente/moderne Methoden (arrow, duckdb, polars)
\item
  Klassische Methoden

  \begin{itemize}
  \tightlist
  \item
    base-r zu vermeiden
  \item
    readr

    \begin{itemize}
    \tightlist
    \item
      in tidyverse
    \item
      nutzerfreundlich, einfach zu lernen / nutzen
    \end{itemize}
  \item
    data.table

    \begin{itemize}
    \tightlist
    \item
      benchmarking:
      https://cran.r-project.org/web/packages/data.table/vignettes/datatable-benchmarking.html
    \end{itemize}
  \item
    datenformate in R: qs, rds
  \item
    dtplyr?
  \end{itemize}
\end{itemize}

\subsubsection{Hoch-performante
Datenverarbeitungsmethoden}\label{hoch-performante-datenverarbeitungsmethoden}

Als Desktopanwendung konzipiert ist R nicht für hoch-performante
Auswertungen von Big Data bekannt. Wie unsere Untersuchung zeigen wird,
beansprucht R nativ übermäßige Systemressourcen, um die
Nutzerfreundlichkeit zu erhöhen. Dadurch dass R Open Source ist ändern
dies jedoch diverse Entwicklerteams weltweit und stellen R
Zusatzfunktionen in Form von R-Paketen zur Verfügung die in R
hoch-performante Auswertung großer Datenmengen ermöglichen.

\begin{itemize}
\tightlist
\item
  R darüber hinaus vermehrt hoch-performante Analysetools, die sich
  ständig verbessern
\item
  ermöglichen sehr schnelle Verarbeitung von Daten, mit einem sehr hohen
  Gesamtvolumen, das sogar den verfügbaren RAM übersteigt.
\item
  sollte möglichst leicht nutzbar sein, nativ in R ohne viele
  Abhängkeiten (organisatorische Einschränkungen)
\end{itemize}

\subsection{Umsetzung im StBA und FDZ}\label{umsetzung-im-stba-und-fdz}

\begin{itemize}
\item
  möglichst leicht nutzbar und gewohnt/verbreitet -\textgreater{}
  R+dpylr

  \begin{itemize}
  \tightlist
  \item
    wechselnde Nutzende (Forscher) unterschiedlicher Fachbereiche mit
    unterschiedlichen Analytischen und technischen Kenntnissen
  \end{itemize}
\item
  stabil
\item
  schnell, optimale Nutzung vorhandener GWAP Zeit und KDFV Zeitraum
\item
  ressourcensparend (RAM Problematik)
\item
  möglichst leicht nutzbar mit den vorhandenen Systemen und
  organisatorischen Strukturen ergibt Anforderungen

  \begin{itemize}
  \tightlist
  \item
    hohe-performanz bereits mit R Grundlagenwissen
  \item
    minimale Anforderungen an zsl. Software
  \end{itemize}
\item
  NeSt Förderung, Datenprodukterstellung, Bereitstellung der FDZ auch
  technischen Infrastruktur und Knowhow der effizienten Nutzung
\item
  Bereitstellung des OSS seit Nov.~2025
\item
  Vorbereitend Exploration und Beratung in geeigneten Aüßerst
  effizienten Methoden -\textgreater{} Ergebnisse in Echtzeit
\item
  Szenario FDZ GWAP Besuch, statt KDFV, sofort Ergebnisse
\item
  Diese Arbeit
\end{itemize}

\textbf{Technische Lösungsansätze.}

\begin{itemize}
\tightlist
\item
  Modernisierung der amtl. Stat, FDZs -\textgreater{} Zugangswege,
  Kapazitäten, \ldots{}
\item
  Technische Weiterentwicklung, stärkere Hardware, leichtere
  Bereitstellung, virtualisierung
\item
  Schnelle Weiterentwicklung von Open-Source Software, R, div. Packages,
  e.g.~Verweis auf Changelog von Arrow, Polars
\end{itemize}

Ausbreitung von R - gleich die richtige Technologie auswählen und
einsetzen. Aber welche ist die ``richtige''? Was für Unterschiede
(Vorteile / Nachteile) haben die Technologien. Dies will dieser Artikel
im Falle des BTP untersuchen. Kontrollierte Benchmark und weitere
(qualitative) Aspekte, wie Userexperience.

Zielgruppe: alle Datenanalytiker mit rudimentären R-Kenntnisse, aber
breit, FDZ, amtl. Stat,

Wir fokussieren uns auf folgende zentrale Forschungsfragen:

\begin{enumerate}
\def\labelenumi{\arabic{enumi}.}
\item
  \textbf{Methodische Effizienz:} Wie unterscheiden sich verschiedene
  statistische Verfahren hinsichtlich ihrer Rechenzeit und
  Speicheranforderungen bei gleichbleibender Ergebnisqualität?
\item
  \textbf{Skalierbarkeit:} Welche Verfahren zeigen die beste Performance
  bei wachsenden Datenmengen, und wo liegen die kritischen Grenzen?
\item
  \textbf{Praktische Anwendbarkeit:} Welche Infrastrukturanforderungen
  stellen moderne Auswertungsmethoden an Forschungsdatenzentren?
\end{enumerate}

\section{Anwendungsfall
Business-Tax-Panel}\label{anwendungsfall-business-tax-panel}

\subsection{Datengrundlage}\label{datengrundlage}

\emph{Bei den verwendeten Statistiken handelt es sich um Vollerhebungen,
teilweise mit Erfassungsgrenze. Insgesamt liegen 77,8 Mio. Beobachtungen
vor, die auf 16,7 Mio. Einheiten zurückzuführen sind. Der Merkmalsumfang
variiert je nach Statistik zwischen 120 und 1.300 Merkmalen was
insgesamt zu einem Merkmalsumfang von über 2.700 Merkmalen führt.
Hierdurch sind Analyse von Anpassungsreaktionen sowie
Verteilungsanalysen für verschiedene Beobachtungszeiträume möglich.
Insgesamt deckt der Forschungsdatensatz eine Vielzahl an
Analysemöglichkeiten, die für die evidenzbasierte steuerpolitische
Entscheidungen nötig sind, ab.}

\emph{Das Business-Tax-Panel kombiniert die Angaben aus verschiedenen
Ertrags- sowie Umsatzsteuerstatistiken im Quer- und Längsschnitt.
Verfügbar sind Informationen aus den folgenden Statistiken :
Gewerbesteuer, Körperschaftsteuer, Umsatzsteuer-Voranmeldung und
-Veranlagung, Statistik über die Personengesellschaften und
Gemeinschaften, Einnahmenüberschussrechnung und einige Merkmale aus dem
Statistischen Unternehmensregister. Alle verwendeten Steuerarte sind ab
dem Berichtsjahr 2013 jährlich verfügbar, dies bildet den Beginn des
Beobachtungszeitraums. Aufgrund der Festsetzungsfrist von vier Jahren
bei den Veranlagungsstatistiken endet der Beobachtungszeitraum aktuell
mit dem Berichtsjahr 2020.}

\subsection{Auswertungsbedarfe}\label{auswertungsbedarfe}

Breites Spektrum: Datenaufbereitung/-bereitstellung und Auswertung von
Mikrodaten durch die Wissenschaft.

\subsubsection{Deskriptive Eckdaten}\label{deskriptive-eckdaten}

\begin{enumerate}
\def\labelenumi{\arabic{enumi}.}
\tightlist
\item
  deskriptive Eckdaten, für BTP als Panel im Quer- und Längsschnitt:
  Fallzahl pro Statistik und Jahr (s.MDR Produkt S. 9f
  https://www.forschungsdatenzentrum.de/sites/default/files/btp\_2013-2019\_on-site\_mdr-prod.pdf)
\end{enumerate}

Kennzahlen um inhaltlose Testdaten zu generieren werden folgende
Eckdaten benötigt und im nächsten BTP berücksichtigt:

\begin{itemize}
\tightlist
\item
  Befüllung
\item
  Statistische Eckdaten univariate (Quantile),
\item
  ggf. multivariate (gemeinsame Verteilung)
\end{itemize}

Exemplarisch, neues BTP beinhaltet dies. Regelmäßige Updates
beschleunigen.

\subsubsection{Regressionen}\label{regressionen}

\begin{enumerate}
\def\labelenumi{\arabic{enumi}.}
\setcounter{enumi}{1}
\tightlist
\item
  Verlustvorträge\footnote{Vortrag Annette 2024 JK} (s. Vortrag NeSt)
  als Indikator für \ldots{} und Einflussgrößen auf Verlustvorträge, für
  Vermutung von Einfluss der Unternehmensgröße (großes Gesamtvolumen,
  Löhne, Spenden, tätige Personen; Auslandskontrolle oder WZ) auf
  positive Verlustvorträge.
\end{enumerate}

\subsection{Technische
Spezifikationen}\label{technische-spezifikationen}

\emph{Für die vielfältigen Analysewünsche der Forschenden, den
Auswertungsbedarf des Fachbereiches bspw. für Sonderauswertungen und der
Erstellung des Metadatenreports des Forschungsdatenzentrums ist der
Forschungsdatensatz konzipiert. Für die Abdeckung aller Anforderungen
ist der Datensatz möglichst wenig beschränkt worden und dadurch
ressourcenintensive. Der damit einhergehende hohe Bedarf an
Speicherkapazität, Arbeitsspeicher und Prozessorleistung stellt sowohl
im Fachbereich als auch im Forschungsdatenzentrum eine Herausforderung
dar.}

\emph{Bei den verwendeten Statistiken handelt es sich um Vollerhebungen,
teilweise mit Erfassungsgrenze. Insgesamt liegen 77,8 Mio. Beobachtungen
vor, die auf 16,7 Mio. Einheiten zurückzuführen sind. Der Merkmalsumfang
variiert je nach Statistik zwischen 120 und 1.300 Merkmalen was
insgesamt zu einem Merkmalsumfang von über 2.700 Merkmalen führt.
Hierdurch sind Analyse von Anpassungsreaktionen sowie
Verteilungsanalysen für verschiedene Beobachtungszeiträume möglich.
Insgesamt deckt der Forschungsdatensatz eine Vielzahl an
Analysemöglichkeiten, die für die evidenzbasierte steuerpolitische
Entscheidungen nötig sind, ab.}

\emph{Für die vielfältigen Analysewünsche der Forschenden, den
Auswertungsbedarf des Fachbereiches bspw. für Sonderauswertungen und der
Erstellung des Metadatenreports des Forschungsdatenzentrums ist der
Forschungsdatensatz konzipiert. Für die Abdeckung aller Anforderungen
ist der Datensatz möglichst wenig beschränkt worden und dadurch
ressourcenintensive. Der damit einhergehende hohe Bedarf an
Speicherkapazität, Arbeitsspeicher und Prozessorleistung stellt sowohl
im Fachbereich als auch im Forschungsdatenzentrum eine Herausforderung
dar.}

\begin{itemize}
\item
  Technische Herausforderung:

  \begin{itemize}
  \tightlist
  \item
    effiziente und verständliche Speicherung
  \item
    Herausfiltern großer Datenmengen irrelevanter Merkmale
  \end{itemize}
\item
  Umfang über verfügbaren Arbeitsspeicher
\item
  Besonders breit mit 3000 Merkmalen, 69 Mio. Beobachtung, sodass
  zeilenbasierte Formate problematisch sind.
\item
  Technische Herausforderung:

  \begin{itemize}
  \tightlist
  \item
    effiziente und verständliche Speicherung
  \item
    Herausfiltern großer Datenmengen irrelevanter Merkmale
  \end{itemize}
\end{itemize}

\section{Infrastruktur des FDZ-Bund}\label{sec-infrastruktur}

Womit wir arbeiten können.

\subsection{Allgemeinn}\label{allgemeinn}

\begin{itemize}
\tightlist
\item
  abgeschottete, innerhalb der eigenen Infrastruktur
\end{itemize}

\subsection{SAS Server}\label{sas-server}

\subsection{StBA R-Server}\label{stba-r-server}

Intern ist R in größeren Kapazitäten nutzbar und Stata nur lokal.

\begin{itemize}
\tightlist
\item
  10 Server á

  \begin{itemize}
  \tightlist
  \item
    1 TB RAM
  \item
    96 Kerne
  \item
    2 GPUs mit insg. 96 GB VRAM
  \end{itemize}
\end{itemize}

Bereitstellung von Paketen und organisatorische einordnung. Veränderung
der Infrastruktur langwierig, Beschaffung via Beschaffungsamt etc.
-\textgreater{} wir müssen für die Zukunft planen und umsetzen.

\subsection{FDZ OnSite-Server}\label{fdz-onsite-server}

StBA R-Server mit noch größeren Kapazitäten und grafischen
Rechenbeschleunigern.

Neu seit November 2025 steht der sog. OnSite-Server im FDZ Bund allen
OnSite Nutzungen bereit -- jeweisl in Stata un R via KDFV und an den
GWAP. Zugriff zu Analyseservern in umfangreichen Kapazitäten.

In FDZ Bund Für Stata (3 Server) oder R (6 Server) in KDFV oder GWAP
stehen jeder Nutzung zur Verfügung:

::: \textbar{} Komponente \textbar{} Umfang \textbar{}
\textbar---------------------\textbar---------------\textbar{}
\textbar{} Arbeitsspeicher \textbar{} bis zu 512 GB \textbar{}
\textbar{} Kernanzahl \textbar{} 16 \textbar{} \textbar{}
CPU-Geschwindigkeit \textbar{} 3.1 GHz \textbar{} :::

\begin{itemize}
\tightlist
\item
  schnelle Direktverbindung zu Netzwerkspeicher
\item
  Fußnote RAM in R 128 GB kann erweitert werden.
\end{itemize}

\section{Datenverarbeitungsmethoden}\label{sec-methoden}

Wir betrachten in diesem Artikel die Performanz sog.
\emph{Datenverarbeitungsmethoden}, d.h. Softwarelösungen, die Daten
einlesen, modifzierern, aggregieren, auswerten oder analysieren. Das
Schreiben von Daten wird in unserer Untersuchung ausgeklammert, da es
sich bei den Veröffentlichungen des Statistischen Bundesamtes oder den
Outputs der Wissenschaft im FDZ um kleine Tabellen handelt, die i.d.R.
keine große Performanz erfordern.{[}\^{}writedata Wir vernachlässigen
Spark, Datenbanken und weitere Technologien, die zusätzliche Einrichtung
erfordern und nicht direkt als R-Package genutzt werden können.

\subsection{Anforderungen an
Datenverarbeitungsmethoden}\label{anforderungen-an-datenverarbeitungsmethoden}

\begin{itemize}
\tightlist
\item
  Performanz und Nutzererfahrung auf dem aktuellen Stand der Technik
\item
  leicht einzurichten auch in den abgeschotteten System der amtl.
  Statistik (tlw. langwierig und schwierig), d.h. wenig Abhängigkeiten
  und unabhängig von Internet bzw. Cloud betreibbar.
\item
  langfristig stabil und gut unterstützt
\item
  leicht nutzbar, gut dokumentiert, gutes Schulungsangebot
\end{itemize}

Die Analyse großer Datensätze wie des in der FDZ-Umgebung stellt
besondere Anforderungen an die verwendeten Werkzeuge. Traditionelle
Ansätze (CSV-Dateien, Base R) stoßen bei dieser Datengröße an ihre
Grenzen. Moderne spaltenorientierte Technologien lösen dieses Problem
indem sie erlauben, nur relevante Variablen bei Bedarf zu einzulesen.

\subsection{Sachstand}\label{sachstand}

\subsubsection{SAS und Stata}\label{sas-und-stata}

\begin{itemize}
\tightlist
\item
  Hauptanalysewerkzeuge sind bisher:

  \begin{itemize}
  \tightlist
  \item
    im StBA SAS 9.4 M8, EG
  \item
    im FDZ zsl. Stata
  \end{itemize}
\item
  Stata

  \begin{itemize}
  \tightlist
  \item
    bisher: V.15 MP-12 (KDFV) und Stata 19 SE lokal (GWAP).
  \item
    neu im OSS in KDFV/GWAP V. 19 gleichermaßen in großen
    Systemkapazitäten.
  \end{itemize}
\end{itemize}

\subsubsection{Empfehlung des FDSZ der deutschen
Bundesbank}\label{empfehlung-des-fdsz-der-deutschen-bundesbank}

\begin{enumerate}
\def\labelenumi{(\arabic{enumi})}
\tightlist
\item
  hat den Spezialfall von in-memory Auswertung betrachtet mit den
  folgenden Ergebnissen:
\end{enumerate}

\begin{itemize}
\tightlist
\item
  R zeigt große Performanzevorteile gegenüber Stata für FDZ Aufgaben:
  ``Note that using parquet and R reduces the computing time compared to
  the optimised Stata approach by 65.1\% from 77.6 to 27.1 seconds''
\item
  zu betrachten ist die Gesamtzeit (d.h. Rechenzeit +
  \textbf{Entwicklungszeit})

  \begin{itemize}
  \tightlist
  \item
    vorhandene Erfahrung und Knowhow (bspw. in Stata) relativiert die
    Performanzgewinne in R. Zu beachten sind hohe Debugging Aufwände an
    den GWAPs der FDZ (kein Internetzugriff)
  \item
    I.d.R. muss in der Entwicklung Code mehrfach ausgeführt werden,
    sodass sich die Laufzeit multipliziert (Wickham\&Grolemund 2016).
  \end{itemize}
\item
  Die u.g. Parquet-Dateien werden bereits standardmäßig im FDSZ
  angeboten.
\item
  Variablenselektion hervorgehoben. Bot Beschleuniging von 39.9\% in R
  und 33.9\% in Stata
\end{itemize}

\subsection{Analysewerkzeug R}\label{analysewerkzeug-r}

Der Einsatz von R wächst stetig mit aktuell über 200 aktiven Nutzenden
im StBA und FDZ.

\begin{itemize}
\tightlist
\item
  R ist die native Programmiersprache von Statistiker/-innen und
  Daten-affinen Programmierer/-innen. Frei verfügbar und leicht
  einzurichten. Somit weit verbreitet, insb. an Universitäten.
\item
  R ist Open Source und wird ständig weiterentwickelt (aktuell in
  Version 4.5 verwendet, jährliche Major Releases).
\item
  Eine weltweit aktive Community löst Probleme zeitnah und ergänzt
  stetig Funktionalitäten in Form von sog. \emph{R-Packages}.
\item
  Bietet nativ bereits viele Datenverarbeitungsmethoden, die jedoch
  durch diese R-Packages ganz neue Level an Nutzerfreundlichkeit,
  Funktionsumfang und Performanz erreichen.
\end{itemize}

\subsection{Klassische Methoden in R}\label{klassische-methoden-in-r}

\subsubsection{Base R}\label{base-r}

\begin{itemize}
\tightlist
\item
  Standardmäßige Basisfunktionen in R.
\item
  Grundlage/einstieg, weder sonderlich amschaulich nocb sonderlich
  performant -\textgreater{} Packages überlassen
\item
  \textbf{Vorteile:} Einfach, weit verbreitet, keine zusätzlichen
  Abhängigkeiten
\item
  \textbf{Nachteile:} Langsam bei großen Daten, hoher Speicherverbrauch,
  alle Daten müssen in RAM passen
\item
  \textbf{Empfehlung:} Zweckdienlich für Einsteiger oder
  gelegenheitsnutzer für ganz kleine Anwendungenoder falls gar keine
  Abhängigkeiten in Form von Pkg erlaubt sind. Ansonsten\ldots{}
\end{itemize}

\subsubsection{Data.table}\label{data.table}

\begin{itemize}
\tightlist
\item
  \{data.table\} in den meisten Fällen die performanteste klassische
  DVM. An die Syntax an base-r angesiedelt und nur hier verwendet,
  spezielles Wissen notwendig.
\item
  \textbf{Vorteile:} Sehr schnell für In-Memory-Operationen, ausgereift,
  kompakte Syntax In teilen der Community sehr beliebt
\item
  \textbf{Nachteile:} Alle Daten müssen in RAM passen, steile Lernkurve,
  keine Out-of-Core-Unterstützung.
\item
  \textbf{Empfehlung:} Gute Alternative wenn Datensatz in RAM passt und
  keine Parquet-Integration benötigt wird
\end{itemize}

\subsubsection{Tidyverse}\label{tidyverse}

Tidyverse ist ein Ökosystem von Packages zur Datenverarbeitung, das für
besonders hohe Lesbarkeit des Codes bekannt ist. Dazu dient eine eigenen
Programmiersyntax, sog. \emph{\{dplyr\}-Syntax}, die den Anspruch hat
menschlicher Sprache abzubilden. Das R-Package \{dplyr\} ist die
Datenverarbeitungsgrammatik innerhalb Tidyverse, die anhand eines
Beispiels in der folgenden Tabelle veranschaulicht wird

{[} Tabelle Beispiel Dplyr {]}

Diese Syntax ist bei einem Teil der R-Community so beliebt, dass sie
vielen anderen Packages ergänzt wurde. Bspw. für \{data.table\} setzen
dies die Packages \{dtplyr\} oder \{tidytable\} um Wie im Falle von
\{dtplyr\} kann die Übersetzung zu \{dplyr\} jedoch die Performanz
verschlechtern.\footnote{S.
  https://www.r-bloggers.com/2024/04/the-truth-about-tidy-wrappers-3/}

\subsection{Best Practices in R}\label{best-practices-in-r}

Ab mittleren Datenmengen (relativ zu CPU und RAM) sind ausgewählte
\textbf{Best Practices} in R:

\begin{itemize}
\tightlist
\item
  \emph{Variablenselektion}: Beim Einlesen von DatenAuswahl der
  benötigten Variablen in der Einlesefunktion ist wesentliche
  Best-Practices (besonders für klassiche Methoden, da dies das Risiko
  für RAM-Abbrüche und Wartezeiten deutlich reduziert). Das Einlesen
  zunächst aller Daten (ggf. mit Variablenselektion zu einem späteren
  Zeitpunkt) verursacht unnötige Wartezeiten, Ressourcenbelegung und
  kann zu Abbrüchen führen.
\item
  \emph{Datenaufteilung}: Wenn der Arbeitsspeicher nicht ausreicht, um
  die benötigten Daten einzulesen, bietet eine Aufteilung auf mehrere
  RAM-Gerechte Datenpakete helfen. Für klassische Methoden kombiniert
  mit der mehrfachen Variablenselektion. Die folgenden hoch-performanten
  Techniken, kommen auch direkt mehreren Dateien nutzen, ohne manuell
  die Datensätze zusammenzuführen.
\end{itemize}

\subsection{Hoch-performante Methoden in R}\label{sec-highperf}

Wir sprechen in diesem Artikel vereinfacht von einer
\emph{hoch-performanten Methode}, sobald die betrachtete
Datenverarbeitungsmethode in der Lage ist, Datenmengen überhalb des
verfügbaren Arbeitsspeichers flexibel und effizient zu nutzen. SAS
erfüllt den Großteil dieser Definition -- lediglich die Effizienz lässt
sich deutlich steigern.

Um Datenauswertung auf höhere Datenmengen skalieren zu können ist ein
Paradigmenwechsel von der ``\emph{vollständigen Verarbeitung im RAM''}
zu der \emph{Auslagerung aus dem RAM} notwendig. Vollständig im RAM sind
Daten zwar sehr schnell und einfach zu nutzen, jedoch beansprucht der
Transport in den RAM (Einlesen) einen Hauptteil der Laufzeit und der
verfügbare RAM reicht i.d.R. nicht für größere Datenmengen. Die
nachfolgenden Technologien sind auch in der Lage Daten überhalb des RAMs
zu verarbeiten. Durch moderne Hardware und Software kann dies sogar sehr
schnell erfolgen.

\subsubsection{Apache Parquet}\label{apache-parquet}

\begin{itemize}
\tightlist
\item
  Datenformat, spaltenbasiert, komprimiert, \ldots{}
\item
  tba - s. Vortrag
\end{itemize}

\subsubsection{Apache Arrow}\label{apache-arrow}

\begin{itemize}
\tightlist
\item
  R-package, C++ basiert, erlaubt Parallelisierung, Datenaufteilung,
  Lazy-Computations, Datenformat im RAM passend zu Parquet
\item
  tba - s. Vortrag
\end{itemize}

\subsubsection{DuckDB}\label{duckdb}

\begin{itemize}
\tightlist
\item
  R-Package, in-memory DB, \ldots{}
\item
  tba - s. Vortrag
\end{itemize}

\subsubsection{Polars}\label{polars}

\begin{itemize}
\tightlist
\item
  an Python-angelehnte Syntax
\item
  tba
\item
  Benchmark des Polars Entwicklerteams \footnote{https://ddotta.github.io/cookbook-rpolars/benchmarking.html}
\end{itemize}

\subsection{Package-Versionen}\label{package-versionen}

Die auf den Servern des StBA und FDZ verfügbaren Packageversionen lassen
sich aus organisatorischen Gründen nicht flexibel ändern, da die
Genehmigung und Einrichtung neuer Packages nur auf Antrag durch unser
R-Betriebsteam erfolgt (i.d.R. innerhalb von 1-2 Tagen ). Die
Package-Versionen, die zum Zeitpunkt des in \textbf{?@sec-benchmark}
dargestellten Performanz-Vergleiches verfügab waren, können Tabelle
Tabelle~\ref{tbl-pkg-versions} entnommen werden.

\begin{table*}

\caption{\label{tbl-pkg-versions}Versionen verfügbarer R-Packages}

\centering{

\centering\begingroup\fontsize{10}{12}\selectfont


\begin{tabular}{|>{}l|||>{}l|}
\hline
\cellcolor[HTML]{dce6f0}{\textcolor[HTML]{404040}{\textbf{Paket}}} & \cellcolor[HTML]{dce6f0}{\textcolor[HTML]{404040}{\textbf{Version}}}\\
\hline
\cellcolor{white}{\textcolor{black}{\{data.table\}}} & \cellcolor[HTML]{f0f8ff}{1.17.8}\\
\cellcolor{white}{\textcolor{black}{\{tidyverse\}}} & \cellcolor[HTML]{f0f8ff}{tba}\\
\cellcolor{white}{\textcolor{black}{\{vroom\}}} & \cellcolor[HTML]{f0f8ff}{tba}\\
\cellcolor{white}{\textcolor{black}{\{arrow\}}} & \cellcolor[HTML]{f0f8ff}{20.0.0.2}\\
\cellcolor{white}{\textcolor{black}{\{duckdb\}}} & \cellcolor[HTML]{f0f8ff}{1.3.2}\\
\cellcolor{white}{\textcolor{black}{\{polars\}}} & \cellcolor[HTML]{f0f8ff}{1.0.1}\\
\hline
\end{tabular}
\endgroup{}

}

\end{table*}%

\section{Performanzvergleich}\label{sec-performance}

\subsection{Versuchsaufbau}\label{versuchsaufbau}

Um zwischen den in Abschnitt Kapitel~\ref{sec-methoden} vorgestellten
Datenverarbeitungsmethoden abzuwägen, haben wir deren Performanz in
einem systematischen Vergleich der Performance auf dem SAS-Server des
Statistischen Bundesamtes und den neuen OnSite R- und Stata-Servern des
FDZ gemessen. Weiterführende Performanzmessung sowie der vollständig
reproduzierbare Code sind auf Opencode verfügbar\footnote{www.opencode.de/URL-To-Be-Announced.}.

Für die Untersuchung wurden simulierte Einzeldaten in verschiedenen
Größen (0,1\%/1\%/10\% von den insgesamt 80 Mio. Beobachtungen des BTP)
betrachtet. Sie wurden aus öffentlich verfügbaren Informationen
generiert, um die wesentlichen Strukturen des BTP für unsere Analyse
künstlich nachzubilden. Abgrenzend zu den im FDZ bereitgestellten
\emph{Datenstrukturfiles} wurden die Testdaten nicht aus echten
Mikrodaten abgeleitet. Der simulierte Datensatz ist ausschließlich für
technische Testzwecke der nachfolgenden Anwendungsfälle bestimmt:

\begin{enumerate}
\def\labelenumi{\arabic{enumi}.}
\tightlist
\item
  \textbf{Fallzahlen} (pro Statistik und Jahr),
\item
  \textbf{Regression} (Einflussfaktoren auf Verlustvorträge).
\end{enumerate}

Um eine zuverlässige Performanz-Messung zu erhalten, wurde jede
Datenverarbeitungsmethode dreifach angewandt. Der Vergleich erfolgt
anhand der mittleren Messung der folgenden Zielgrößen:

\begin{enumerate}
\def\labelenumi{\arabic{enumi}.}
\tightlist
\item
  \textbf{Arbeitsspeicher} der Auswertung,
\item
  \textbf{Laufzeit} der Auswertung (real verstrichene Zeit),
\item
  \textbf{Festplattenspeicher} der verwendeten Daten.
\end{enumerate}

Die Zielgrößen sind in dieser Reihenfolge zu priorisieren, da ein
übermäßiger Arbeitsspeicherverbrauch zu schwerwiegenden Abbrüchen führen
kann, während hohe Laufzeiten dennoch zu einem erfolgreichen Ergebnis
führen können. Festplattenspeicher ist ausreichend vorhanden, aber
unnötige Belegung sollte vermieden werden.

Feine Abstufungen der Zielgrößen sind für unsere Aufgaben nicht
gleichermaßen relevant. Eine Methode, die die betrachteten Auswertungen
in unter 10 Sekunden mit weniger als 4 GB Arbeitsspeicher durchführen
kann, wird bereits als optimal bewertet. Weitere Optimierungen sind in
diesem Fall nicht erforderlich. Falls bereits die gewünschte Performanz
erreicht wurde, rückt die Reduktion des Entwicklungsaufwand in den
Vordergrund. Entsprechend sind Aspekte wie einfache Anwendbarkeit und
Nachvollziehbarkeit -- selbst für unerfahrene Nutzende --
ausschlaggebend.

\subsection{Vergleichswert in SAS und
Stata}\label{vergleichswert-in-sas-und-stata}

Klassisch werden die betrachteten Auswertungsaufgaben im Statischen
Bundesamt hauptsächlich mit SAS und im FDZ mit Stata gelöst. Um die
nachfolgenden Ergebnisse einzuordnen, haben wir die Performanzmessungen
für 1\% des simulierten BTP-Daten (etwa. 800.000 Beobachtungen) in SAS
und Stata durchgeführt.

\begin{table*}

\caption{\label{tbl-baseline}Baselines in SAS und Stata (anhand 1\% des simuliertes BTP)}

\centering{

\centering\begingroup\fontsize{8}{10}\selectfont


\begin{tabular}{|>{}l||>{}l||>{}r||>{}r||>{}r|}
\toprule
Software & Anwendung & Arbeitsspeicher & Laufzeit & Festplattenspeicher\\
\midrule
 & \cellcolor[HTML]{f0f8ff}{Fallzahl} & \cellcolor[HTML]{f0f8ff}{65 MB} & \cellcolor[HTML]{f0f8ff}{2.4 Min.} & \cellcolor[HTML]{f0f8ff}{33 GB}\\
\cmidrule{2-5}
\multirow[t]{-2}{*}{\raggedright\arraybackslash \cellcolor{white}{\textcolor{black}{SAS}}} & \cellcolor[HTML]{f0f8ff}{Regression} & \cellcolor[HTML]{f0f8ff}{7 MB} & \cellcolor[HTML]{f0f8ff}{2.8 Min.} & \cellcolor[HTML]{f0f8ff}{33 GB}\\
\cmidrule{1-5}
 & \cellcolor[HTML]{f0f8ff}{Fallzahl} & \cellcolor[HTML]{f0f8ff}{33 GB} & \cellcolor[HTML]{f0f8ff}{55,2 Sek.} & \cellcolor[HTML]{f0f8ff}{32 GB}\\
\cmidrule{2-5}
\multirow[t]{-2}{*}{\raggedright\arraybackslash \cellcolor{white}{\textcolor{black}{Stata}}} & \cellcolor[HTML]{f0f8ff}{Regression} & \cellcolor[HTML]{f0f8ff}{33 GB} & \cellcolor[HTML]{f0f8ff}{55,8 Sek.} & \cellcolor[HTML]{f0f8ff}{32 GB}\\
\bottomrule
\end{tabular}
\endgroup{}

}

\end{table*}%

Die Ergebnisse in Tabelle~\ref{tbl-baseline} zeigen, dass der
umfangreiche Arbeitsspeicher des FDZ OnSite-Servers nicht ausreicht, um
das ganze BTP in Stata einzulesen, sodass zwingend mit
Variablenselektion gearbeitet werden muss. SAS zeigt seine Stärke anhand
der geringen Beanspruchung des Arbeitsspeichers, benötigt aber mit etwa
2,3 Minuten pro Auswertung die doppelte Laufzeit gegenüber Stata.

\subsection{Kleiner Vergleich zwischen R, SAS und
Stata}\label{kleiner-vergleich-zwischen-r-sas-und-stata}

Für 1\% der simulierten BTP Daten, reduzieren die
R-Datenverarbeitungsmethode alle Zielgrößen gegenüber SAS und Stata
erheblich (siehe Tabelle~\ref{tbl-benchmark-small} und
Tabelle~\ref{tbl-benchmark-regr-small}). Selbst klassische R-Methoden
sind um den Faktor 7 schneller und benötigen nur 33 \% des
Festplattenspeichers. \{polars\}, \{duckdb\} und \{arrow\} lösen die
betrachteten Auswertungen sogar

\begin{itemize}
\tightlist
\item
  in \textbf{unter 2 Sekunden} (27-fach schneller),
\item
  mit unter \textbf{300 MB Arbeitsspeicher} -- \{duckdb\} und
  \{polars\}{[}\^{}polars-ram{]} sogar mit unter \textbf{5 MB} (10-fach
  weniger)\footnote{Lediglich in der Arbeitsspeicherbeanspruchung, ist
    SAS deutlich besser als klassische R-Methoden. Gegenüber
    hoch-performanten R-Methoden benötigt SAS jedoch den 10-fachen
    Arbeitsspeicher.}\footnote{\{arrow\} benötigt mit 274 MB deutlich
    mehr Arbeitsspeicher als \{polars\} und \{duckdb\}, aber kommt
    dennoch mit minimalen Entwicklungsumgebungen aus, bspw. einem
    standardmäßigen Laptop.} --
\item
  und nur mit \textbf{5 GB Festplattenspeicher} (6-fach weniger
  gegenüber SAS und Stata).
\end{itemize}

Unter den klassischen R-Methoden bietet \{data.table\} die beste
Performanz. Dabei ist wesentlich zu beachten, dass die Einlesefunktion
\texttt{fread} direkt die relevanten Variablen auswählt. Dadurch lassen
sich unsere Auswertungen um einen Faktor 3 beschleunigen und vor allem
der Arbeitsspeicher um einen Faktor 1.000 reduzieren.

Alle hoch-performanten Methoden sind dermaßen effizient, dass die
Datenmenge von 800.000 Beobachtungen nicht mehr ausreicht, um zwischen
den besten Methoden zu differenzieren. Daher erhöhen wir im nächsten
Experiment die Datenmenge auf 8.000.000 Beobachtungen (10\% BTP).

\begin{table*}

\caption{\label{tbl-benchmark-small}Performanz Messungen für die Fallzahlberechnung anhand 1\% des BTP}

\centering{

\centering\begingroup\fontsize{8}{10}\selectfont


\begin{tabular}{|>{}l|||>{}l|||>{}l|||>{}r|||>{}r|||>{}r|}
\hline
\cellcolor[HTML]{dce6f0}{\textcolor[HTML]{404040}{\textbf{Package}}} & \cellcolor[HTML]{dce6f0}{\textcolor[HTML]{404040}{\textbf{Dateiformat}}} & \cellcolor[HTML]{dce6f0}{\textcolor[HTML]{404040}{\textbf{Optimierung}}} & \cellcolor[HTML]{dce6f0}{\textcolor[HTML]{404040}{\textbf{Arbeitsspeicher}}} & \cellcolor[HTML]{dce6f0}{\textcolor[HTML]{404040}{\textbf{Laufzeit}}} & \cellcolor[HTML]{dce6f0}{\textcolor[HTML]{404040}{\textbf{Festplattenspeicher}}}\\
\hline
\cellcolor{white}{\textcolor{black}{\{polars\}}} & \cellcolor[HTML]{f0f8ff}{Parquet} & \cellcolor[HTML]{f0f8ff}{-} & \cellcolor[HTML]{f0f8ff}{3,3 MB} & \cellcolor[HTML]{f0f8ff}{1,0 Sek} & \cellcolor[HTML]{f0f8ff}{4,9 GB}\\
\cellcolor{white}{\textcolor{black}{\{duckdb\}}} & \cellcolor[HTML]{f0f8ff}{Parquet} & \cellcolor[HTML]{f0f8ff}{-} & \cellcolor[HTML]{f0f8ff}{4,8 MB} & \cellcolor[HTML]{f0f8ff}{1,5 Sek} & \cellcolor[HTML]{f0f8ff}{4,9 GB}\\
\cellcolor{white}{\textcolor{black}{\{arrow\}}} & \cellcolor[HTML]{f0f8ff}{Parquet} & \cellcolor[HTML]{f0f8ff}{Variablenauswahl} & \cellcolor[HTML]{f0f8ff}{273,7 MB} & \cellcolor[HTML]{f0f8ff}{1,7 Sek} & \cellcolor[HTML]{f0f8ff}{4,9 GB}\\
\cellcolor{white}{\textcolor{black}{\{data.table\}}} & \cellcolor[HTML]{f0f8ff}{CSV} & \cellcolor[HTML]{f0f8ff}{Variablenauswahl} & \cellcolor[HTML]{f0f8ff}{42,9 MB} & \cellcolor[HTML]{f0f8ff}{2,5 Sek} & \cellcolor[HTML]{f0f8ff}{13,0 GB}\\
\cellcolor{white}{\textcolor{black}{\{arrow\}}} & \cellcolor[HTML]{f0f8ff}{Parquet} & \cellcolor[HTML]{f0f8ff}{Datenaufteilung (gleichmäßig)} & \cellcolor[HTML]{f0f8ff}{229,7 MB} & \cellcolor[HTML]{f0f8ff}{2,8 Sek} & \cellcolor[HTML]{f0f8ff}{4,9 GB}\\
\cellcolor{white}{\textcolor{black}{\{arrow\}}} & \cellcolor[HTML]{f0f8ff}{Parquet} & \cellcolor[HTML]{f0f8ff}{-} & \cellcolor[HTML]{f0f8ff}{272,8 MB} & \cellcolor[HTML]{f0f8ff}{2,8 Sek} & \cellcolor[HTML]{f0f8ff}{4,9 GB}\\
\cellcolor{white}{\textcolor{black}{\{arrow\}}} & \cellcolor[HTML]{f0f8ff}{Parquet} & \cellcolor[HTML]{f0f8ff}{-} & \cellcolor[HTML]{f0f8ff}{230,4 MB} & \cellcolor[HTML]{f0f8ff}{3,8 Sek} & \cellcolor[HTML]{f0f8ff}{4,9 GB}\\
\cellcolor{white}{\textcolor{black}{\{arrow\}}} & \cellcolor[HTML]{f0f8ff}{Parquet} & \cellcolor[HTML]{f0f8ff}{Datenaufteilung (Berichtsjahr)} & \cellcolor[HTML]{f0f8ff}{229,9 MB} & \cellcolor[HTML]{f0f8ff}{3,9 Sek} & \cellcolor[HTML]{f0f8ff}{4,8 GB}\\
\cellcolor{white}{\textcolor{black}{\{data.table\}}} & \cellcolor[HTML]{f0f8ff}{CSV} & \cellcolor[HTML]{f0f8ff}{-} & \cellcolor[HTML]{f0f8ff}{33,7 GB} & \cellcolor[HTML]{f0f8ff}{7,4 Sek} & \cellcolor[HTML]{f0f8ff}{13,0 GB}\\
\cellcolor{white}{\textcolor{black}{\{data.table\}}} & \cellcolor[HTML]{f0f8ff}{CSV} & \cellcolor[HTML]{f0f8ff}{Datenaufteilung (Berichtsjahr)} & \cellcolor[HTML]{f0f8ff}{50,4 GB} & \cellcolor[HTML]{f0f8ff}{13,6 Sek} & \cellcolor[HTML]{f0f8ff}{10,5 GB}\\
\hline
\end{tabular}
\endgroup{}

}

\end{table*}%

\begin{table*}

\caption{\label{tbl-benchmark-regr-small}Performanz Messungen für die Regressionsberechnung anhand 1\% des BTP}

\centering{

\centering\begingroup\fontsize{8}{10}\selectfont


\begin{tabular}{|>{}l|||>{}l|||>{}l|||>{}r|||>{}r|||>{}r|}
\hline
\cellcolor[HTML]{dce6f0}{\textcolor[HTML]{404040}{\textbf{Package}}} & \cellcolor[HTML]{dce6f0}{\textcolor[HTML]{404040}{\textbf{Dateiformat}}} & \cellcolor[HTML]{dce6f0}{\textcolor[HTML]{404040}{\textbf{Optimierung}}} & \cellcolor[HTML]{dce6f0}{\textcolor[HTML]{404040}{\textbf{Arbeitsspeicher}}} & \cellcolor[HTML]{dce6f0}{\textcolor[HTML]{404040}{\textbf{Laufzeit}}} & \cellcolor[HTML]{dce6f0}{\textcolor[HTML]{404040}{\textbf{Festplattenspeicher}}}\\
\hline
\cellcolor{white}{\textcolor{black}{\{arrow\}}} & \cellcolor[HTML]{f0f8ff}{Parquet} & \cellcolor[HTML]{f0f8ff}{Datenaufteilung (Berichtsjahr)} & \cellcolor[HTML]{f0f8ff}{529,7 MB} & \cellcolor[HTML]{f0f8ff}{3,8 Sek} & \cellcolor[HTML]{f0f8ff}{4,8 GB}\\
\cellcolor{white}{\textcolor{black}{\{dplyr\}}} & \cellcolor[HTML]{f0f8ff}{CSV} & \cellcolor[HTML]{f0f8ff}{Variablenauswahl} & \cellcolor[HTML]{f0f8ff}{154,8 MB} & \cellcolor[HTML]{f0f8ff}{13,2 Sek} & \cellcolor[HTML]{f0f8ff}{13,0 GB}\\
\hline
\end{tabular}
\endgroup{}

}

\end{table*}%

\clearpage

\subsection{Größere Datenmengen in
R}\label{gruxf6uxdfere-datenmengen-in-r}

Auch für 10 \% der simulierten BTP-Daten, kann R die beiden betrachteten
Auswertungen in Echtzeit mit minimalem Arbeitsspeicher beantworten
(siehe Tabelle~\ref{tbl-benchmark} und
Tabelle~\ref{tbl-benchmark-regr}). Die Auswertungsmethoden mit der
höchsten Performanz nutzen alle das \{parquet\} und erreichen:

\begin{enumerate}
\def\labelenumi{\arabic{enumi}.}
\tightlist
\item
  \{arrow\} benötigt nur \textbf{4,7 Sekunden} und 230 MB
  Arbeitsspeicher,
\item
  \{duckdb\} benötigt nur 5,1 Sekunden und \textbf{4,7 MB
  Arbeitsspeicher}.
\end{enumerate}

{[}tba: Plot der Performanz unterschiede{]}

Unerwarteterweise ist die Laufzeit von \{polars\} für diese Datenmenge
mit 1,2 Minuten weit abgeschlagen, obwohl das gleiche Auswertungsskript
für kleinere Datenmengen immer die geringste Laufzeit aufwies. Da
\{polars\} das jüngste der betrachteten R-Packages ist, vermuten wir
Performanzprobleme in der vorliegenden Version 1.0.1\footnote{Das
  R-Package \{polars\} Version 1.0.1 war nur der erste Minor Release
  nach einer kompletten Überarbeitung. Inzwischen sind 5 neuere Major
  Releases mit diversen Performanz-Verbesserungen verfügbar.}. Zudem hat
es noch größere Entwicklungsaufwände beansprucht, da der an Python
angelehnte Syntaxstil deutlich von üblicher Programmierung in R
abweicht. Öffentliche Hilfestellung bezieht sich hauptsächlich auf
Python (oder Rust), was sich in den Antworten der KI-Chatassistenten
wiederspiegelt.

Klassische Methoden kommen bei der betrachteten Datenmenge bereits an
ihre Grenzen. Falls in \{data.table\} unbedacht der vollständige
Datensatz eingelesen wird, werden in unserem Anwemdungsfall die
verfügbaren 512 GB Arbeitsspeicher vollständig beansprucht. Schon wenn
klassische Methoden auf Datenmengen von ca. 25-33\% des verfügbaren
Arbeitsspeichers angewandt werden, ist somit deutlich mehr Erfahrung der
Nutzenden notwendig, um die relevanten Variablen manuell geschickt zu
selektieren unter Beobachtung des Arbeitsspeicherverbrauchs.

Somit kommen wir insgesamt zu dem Ergebnis, dass \{arrow\} kombiniert
mit \{parquet\} für die Bewältigung unserer Aufgaben die beste
Datenverarbeitungsmethode ist, da die Methodik hervorragende Performanz
und zugleich hohe Nutzerfreundlichkeit bietet. Im Detail haben uns der
folgenden Mehrwerte überzeugt:

\begin{enumerate}
\def\labelenumi{\arabic{enumi}.}
\tightlist
\item
  Für die relevanteste Datenmenge von 8 Mio. Beobachtungen zeigte
  \{arrow\} die geringsten Laufzeiten unter 5 Sekunden, was der
  Fachstatistik Antworten auf Sonderanfragen in Echtzeit ermöglicht und
  der Wissenschaft alle denkenswerten explorativen Analysen auszuweiten.
\item
  auch die Nutzererfahrung gefiel uns mit \{arrow\} am besten und somit
  waren die einhergehenden Entwicklungaufwände gegenüber den Anderen
  Methoden am geringsten.
\item
  ,\{arrow\} verwendet die anschauliche \{tidyverse\} Syntax nativ.
  Vorhandener Code lässt sich zum Großteil änderungsfrei und ohne
  Performanzverluste auf \{arrow\} übertragen. \{duckdb\} kann bei
  Kenntnissen in SQL zu bevorzugen sein und \{polars\} mit Erfahrungen
  in Python.
\item
  Beide Anwendungen ließen sich in unter 20 nachvollziehbaren Zeilen in
  \{arrow\} programmieren.\footnote{www.opencode.de/URL-To-Be-Announced.}
\item
  Obwohl nicht alle Features aus \{tidyverse\} für \{arrow\} verfügbar
  sind, ließ sich bisher immer eine Lösung finden, um die gewünschte
  Auswertung umzusetzen, bspw. durch Kombination mit \{duckdb\} oder
  andere Funktionsaufrufe.
\item
  \{arrow\} spart ausreichend Arbeitsspeicher (mit Verbrauch unter 1 GB)
  (tba: Fußnote zu duckdb/polars). Der Arbeitsspeicherverbrauch von
  \{arrow\} ist in unserer Auswertung nahezu konstant abhängig von der
  Datenmenge, sodass \{arrow\} leicht auf noch größere Datenmengen
  skalieren kann
\item
  Dadurch das \{arrow\} die Nutzung auch von mehreren \{parquet\}
  Dateien unterstützt spart gegenüber CSV mind. 50\% der Festplatte (für
  das reale BTP sogar 80\%+) und ermöglicht schnelle Nutzung auch von
  Datenmengen überhalb des verfügbaren Arbeitsspeichers.
\end{enumerate}

\begin{table*}

\caption{\label{tbl-benchmark}Performanz Messungen für die Fallzahlberechnung anhand 10\% des BTP}

\centering{

\centering\begingroup\fontsize{8}{10}\selectfont


\begin{tabular}{|>{}l|||>{}l|||>{}l|||>{}r|||>{}r|||>{}r|}
\hline
\cellcolor[HTML]{dce6f0}{\textcolor[HTML]{404040}{\textbf{Package}}} & \cellcolor[HTML]{dce6f0}{\textcolor[HTML]{404040}{\textbf{Dateiformat}}} & \cellcolor[HTML]{dce6f0}{\textcolor[HTML]{404040}{\textbf{Optimierung}}} & \cellcolor[HTML]{dce6f0}{\textcolor[HTML]{404040}{\textbf{Arbeitsspeicher}}} & \cellcolor[HTML]{dce6f0}{\textcolor[HTML]{404040}{\textbf{Laufzeit}}} & \cellcolor[HTML]{dce6f0}{\textcolor[HTML]{404040}{\textbf{Festplattenspeicher}}}\\
\hline
\cellcolor{white}{\textcolor{black}{\{arrow\}}} & \cellcolor[HTML]{f0f8ff}{Parquet} & \cellcolor[HTML]{f0f8ff}{Datenaufteilung (gleichmäßig)} & \cellcolor[HTML]{f0f8ff}{229,9 MB} & \cellcolor[HTML]{f0f8ff}{4,7 Sek} & \cellcolor[HTML]{f0f8ff}{48,8 GB}\\
\cellcolor{white}{\textcolor{black}{\{duckdb\}}} & \cellcolor[HTML]{f0f8ff}{Parquet} & \cellcolor[HTML]{f0f8ff}{-} & \cellcolor[HTML]{f0f8ff}{4,7 MB} & \cellcolor[HTML]{f0f8ff}{5,1 Sek} & \cellcolor[HTML]{f0f8ff}{49,4 GB}\\
\cellcolor{white}{\textcolor{black}{\{arrow\}}} & \cellcolor[HTML]{f0f8ff}{Parquet} & \cellcolor[HTML]{f0f8ff}{Variablenauswahl} & \cellcolor[HTML]{f0f8ff}{702,1 MB} & \cellcolor[HTML]{f0f8ff}{7,5 Sek} & \cellcolor[HTML]{f0f8ff}{49,4 GB}\\
\cellcolor{white}{\textcolor{black}{\{arrow\}}} & \cellcolor[HTML]{f0f8ff}{Parquet} & \cellcolor[HTML]{f0f8ff}{Datenaufteilung (Berichtsjahr)} & \cellcolor[HTML]{f0f8ff}{4,7 MB} & \cellcolor[HTML]{f0f8ff}{11,7 Sek} & \cellcolor[HTML]{f0f8ff}{48,1 GB}\\
\cellcolor{white}{\textcolor{black}{\{arrow\}}} & \cellcolor[HTML]{f0f8ff}{Parquet} & \cellcolor[HTML]{f0f8ff}{-} & \cellcolor[HTML]{f0f8ff}{230,4 MB} & \cellcolor[HTML]{f0f8ff}{17,0 Sek} & \cellcolor[HTML]{f0f8ff}{49,4 GB}\\
\cellcolor{white}{\textcolor{black}{\{arrow\}}} & \cellcolor[HTML]{f0f8ff}{Parquet} & \cellcolor[HTML]{f0f8ff}{-} & \cellcolor[HTML]{f0f8ff}{701,2 MB} & \cellcolor[HTML]{f0f8ff}{20,6 Sek} & \cellcolor[HTML]{f0f8ff}{49,4 GB}\\
\cellcolor{white}{\textcolor{black}{\{data.table\}}} & \cellcolor[HTML]{f0f8ff}{CSV} & \cellcolor[HTML]{f0f8ff}{Variablenauswahl} & \cellcolor[HTML]{f0f8ff}{427,5 MB} & \cellcolor[HTML]{f0f8ff}{23,6 Sek} & \cellcolor[HTML]{f0f8ff}{130,2 GB}\\
\cellcolor{white}{\textcolor{black}{\{polars\}}} & \cellcolor[HTML]{f0f8ff}{Parquet} & \cellcolor[HTML]{f0f8ff}{-} & \cellcolor[HTML]{f0f8ff}{3,3 MB} & \cellcolor[HTML]{f0f8ff}{1,2 Min} & \cellcolor[HTML]{f0f8ff}{49,4 GB}\\
\cellcolor{white}{\textcolor{black}{\{data.table\}}} & \cellcolor[HTML]{f0f8ff}{CSV} & \cellcolor[HTML]{f0f8ff}{-} & \cellcolor[HTML]{f0f8ff}{336,0 GB} & \cellcolor[HTML]{f0f8ff}{2,2 Min} & \cellcolor[HTML]{f0f8ff}{130,2 GB}\\
\cellcolor{white}{\textcolor{black}{\{data.table\}}} & \cellcolor[HTML]{f0f8ff}{CSV} & \cellcolor[HTML]{f0f8ff}{Datenaufteilung (Berichtsjahr)} & \cellcolor[HTML]{f0f8ff}{506,4 GB} & \cellcolor[HTML]{f0f8ff}{5,0 Min} & \cellcolor[HTML]{f0f8ff}{104,6 GB}\\
\hline
\end{tabular}
\endgroup{}

}

\end{table*}%

\begin{table*}

\caption{\label{tbl-benchmark-regr}Performanz Messungen für die Regressionsberechnung anhand 10\% des BTP}

\centering{

\centering\begingroup\fontsize{8}{10}\selectfont


\begin{tabular}{|>{}l|||>{}l|||>{}l|||>{}r|||>{}r|||>{}r|}
\hline
\cellcolor[HTML]{dce6f0}{\textcolor[HTML]{404040}{\textbf{Package}}} & \cellcolor[HTML]{dce6f0}{\textcolor[HTML]{404040}{\textbf{Dateiformat}}} & \cellcolor[HTML]{dce6f0}{\textcolor[HTML]{404040}{\textbf{Optimierung}}} & \cellcolor[HTML]{dce6f0}{\textcolor[HTML]{404040}{\textbf{Arbeitsspeicher}}} & \cellcolor[HTML]{dce6f0}{\textcolor[HTML]{404040}{\textbf{Laufzeit}}} & \cellcolor[HTML]{dce6f0}{\textcolor[HTML]{404040}{\textbf{Festplattenspeicher}}}\\
\hline
\cellcolor{white}{\textcolor{black}{\{arrow\}}} & \cellcolor[HTML]{f0f8ff}{Parquet} & \cellcolor[HTML]{f0f8ff}{Datenaufteilung (Berichtsjahr)} & \cellcolor[HTML]{f0f8ff}{782,1 MB} & \cellcolor[HTML]{f0f8ff}{5,2 Sek} & \cellcolor[HTML]{f0f8ff}{48,1 GB}\\
\cellcolor{white}{\textcolor{black}{\{dplyr\}}} & \cellcolor[HTML]{f0f8ff}{CSV} & \cellcolor[HTML]{f0f8ff}{Variablenauswahl} & \cellcolor[HTML]{f0f8ff}{1,5 GB} & \cellcolor[HTML]{f0f8ff}{1,9 Min} & \cellcolor[HTML]{f0f8ff}{130,2 GB}\\
\hline
\end{tabular}
\endgroup{}

}

\end{table*}%

{[}tba: Plot RAM-Bedarf abhängig von der Datenmenge{]}

\setcounter{subsection}{4}

\section{Anwendung auf das BTP}\label{sec-ergebnis}

Anwendung der besten Methode auf die echten Daten des vollständige BTP
ergibt:

\begin{itemize}
\tightlist
\item
  Performanz von \{arrow\} Fallzahlen (ein paar Sekunden) und Regression
  (Sekunden)
\item
  leicht lesbar in \{arrow\}
\end{itemize}

\subsection{Ergebnis}\label{ergebnis}

{[}Ergebnistabelle Fallzahl, ggf. als Heatmap{]}

-\textgreater{} Zukünftige Nutzung für MDR wird umgesetzt.

{[}Model Summary?{]}

\subsection{Vergleich zu SAS, (Stata),
Datatable}\label{vergleich-zu-sas-stata-datatable}

\begin{itemize}
\tightlist
\item
  Ergebnis in SAS
\item
  (optional) Datatable ohne sekeltives Einlesen führt zu RAM Crash
  (RAM-Fehler?)
\item
  (optional) Datatable Sel.
\end{itemize}

Real world application unserer Ergebnisse auf das Echte BTP

Ergebnis für gesamten Datensatz

-\textgreater{} Größere Speicherreduktion, da das echte BTP mehr
fehlende Werte enthält, als anhand der externen Daten angenommen.

\setcounter{subsection}{5}

\section{Fazit}\label{sec-ausblick}

\subsection{Kernaussage des Artikels}\label{kernaussage-des-artikels}

Wir empfehlen eine Datenhaltung im Parquet-Format und \textbf{R}
incl.~\{arrow\} + \{parquet\} (bei Bedarf \{duckdb\}, \{polars\}) zu
erproben und zu etablieren! Denn unsere Untersuchung demonstriert ganz
neue Möglichkeiten, analytische Fragestellungen anhand großer
Datenmengen in Echtzeit zu beantworten.

\subsection{Auswirkung auf IT-Infrastruktur und
Beratung}\label{auswirkung-auf-it-infrastruktur-und-beratung}

\begin{itemize}
\tightlist
\item
  Fokus auf Hardware und Software mit größtem Mehrwert (weniger RAM und
  R anstatt Spark, GPUs, Datenbanken, Cloud, \ldots)
\item
  Fokus auf Maintainance und Patching von R, Parquet und hier genannte
  Packages \{arrow\}, \{duckdb\}, \{polars\}
\item
  Etablierung dieser Softwarelösungen kann langfristig\ldots{}

  \begin{itemize}
  \tightlist
  \item
    die Hardwareanforderungen reduzieren
  \item
    die Analysesoftwarevielfalt verschlanken, somit Beschaffungen
    einsparen (besserer Ausdruck?)
  \end{itemize}
\item
  Schulungsangebot benötigt für \emph{``hoch-performante Auswertungen in
  R''}
\item
  eigener Beitrag zu Open Source entscheidend, um die Etablierung zu
  fördern.
\end{itemize}

\subsection{Auswirkungen auf die
Statistikproduktion}\label{auswirkungen-auf-die-statistikproduktion}

\begin{itemize}
\tightlist
\item
  Steigerung der Effizienz und Reaktionsgeschwindigkeit
\item
  (Vorsichtig!? Annette, lieber weg oder rein?) Vereinfachung der
  Statistikproduktion, durch leicht nachvollziehbaren Code und Fokus auf
  Wissensaufbau in R
\item
  Aufbau eines Datenbestands in Parquet
\item
  Anwendung auf weitere Produkte

  \begin{itemize}
  \tightlist
  \item
    Überarbeitung des BTP
  \item
    Mindeststeuer (oder?)
  \item
    DRG
  \item
    TPP 2
  \end{itemize}
\end{itemize}

\subsection{Auswirkungen auf die FDZ}\label{auswirkungen-auf-die-fdz}

\begin{itemize}
\tightlist
\item
  Berücksichtigung in Mustersyntaxen und Nutzendenberatung im FDZ Bund.
  Teilen von Code-Snippets und Templates fördern.
\item
  mögliche Vereinfachung der Betreuung (Fokus auf R, Reduktion von
  Performanz- und Programmierproblemen)
\item
  Bereitstellung von Parquet in geeigneter Aufteilung empfohlen
\item
  Erprobung im FDZ Bund Vorbild für andere FDZ
\item
  Optimierung von Hardware und Software Angebot
\end{itemize}

\subsection{Auswirkung auf die
Wissenschaft}\label{auswirkung-auf-die-wissenschaft}

\begin{itemize}
\tightlist
\item
  Verwendung von R bietet am FDZ Bund einen deutlichen
  \textbf{Wettbewerbsvorteil in der emp. Forschung}. Lediglich R ist (am
  GWAP und via KDFV) am FDZ Bund in der Lage große Datenmengen innerhalb
  von Sekunden zu analysieren und Datenmengen überhalb des verfügbaren
  Arbeitspeicher, wie das BTP, ohne vorherige Variablenselektion
  auszuwerten.

  \begin{itemize}
  \tightlist
  \item
    Vollmaterial kann statt Stichproben angeboten werden. Erleichtert
    die FDZ Betreuung und GWAP/KDFV-Verfügbarkeit
  \item
    mehr explorative oder komplexe Analysen sind in kürzerer Zeit
    möglich.
  \end{itemize}
\item
  Dazu ist R als frei verfügbare Open-Source Software leicht zu lernen
  und \{arrow\} bietet hoch-performante Analysen selbst für
  R-Basiskenntnisse. Empfehlung an die emp. Wissenschaft, sich verstärkt
  mit R zu befassen.
\item
  Für analoge über-RAM-Verarbeitungen bietet Stata aktuell keine Lösung.
  SAS bietet bei weitem nicht hier gemessene Performanz von R. R
  hervorragende Möglichkeiten bietet Performanz und
  Anwenderfreundlichkeit zu übertreffen.
\item
  Weiteres Vorgehen: \textbf{NeSt-Pilotprojekt}. Aus dem realen Use Case
  Optimierung der Systemkonfiguration und Beratungsbedarfe
  identifizieren.
\end{itemize}

\subsection{Schlussbemerkung}\label{schlussbemerkung}

Die Analyse großer Datensätze in der amtlichen Statstik und den
Forschungsdatenzen können von einem signifikanten Performanz-Boost
profitieren, wenn die vorgestellten Werkzeuge adaptiert werden.
\textbf{Arrow, DuckDB, Polars und Parquet} -- bieten das Potenzial, die
Effizienz erheblich zu steigern und neue Analysemöglichkeiten zu
eröffnen.

Der Übergang erfordert jedoch eine \textbf{koordinierte Anstrengung} von
Infrastrukturbetreibern, amtlicher Statistik, Forschungsdatenzentren und
Forschenden. Die systematische Wissensaufbau, Dokumentation und
Verbreitung von Best Practices wird entscheidend für den Erfolg sein.

Die amtliche Statistik kann von diesen Entwicklungen in mehrfacher
Hinsicht profitieren: \textbf{schnellere Erkenntnisgewinnung},
\textbf{umfassendere Analysen} und \textbf{effizientere
Ressourcennutzung}. Dies trägt letztlich zur Stärkung der
evidenzbasierten Politik- und Gesellschaftsberatung bei.

\section*{Anhang}\label{anhang}
\addcontentsline{toc}{section}{Anhang}

Inhalte für die Veröffentlichung in Open Code.

\subsection*{Weitere Performancemessungen für kleinere
Datenmengen.}\label{weitere-performancemessungen-fuxfcr-kleinere-datenmengen.}
\addcontentsline{toc}{subsection}{Weitere Performancemessungen für
kleinere Datenmengen.}

\begin{table*}

\caption{\label{tbl-bm-cnt-small}Benchmarkergebnisse für Fallzahlen (kleine Stichproben des BTP)}

\centering{

\centering\begingroup\fontsize{8}{10}\selectfont


\begin{tabular}{|>{}r||>{}l||>{}l||>{}l||>{}r||>{}r||>{}r|}
\toprule
\cellcolor[HTML]{dce6f0}{\textcolor[HTML]{404040}{\textbf{Beob.}}} & \cellcolor[HTML]{dce6f0}{\textcolor[HTML]{404040}{\textbf{Package}}} & \cellcolor[HTML]{dce6f0}{\textcolor[HTML]{404040}{\textbf{Dateiformat}}} & \cellcolor[HTML]{dce6f0}{\textcolor[HTML]{404040}{\textbf{Optimierung}}} & \cellcolor[HTML]{dce6f0}{\textcolor[HTML]{404040}{\textbf{Arbeitsspeicher}}} & \cellcolor[HTML]{dce6f0}{\textcolor[HTML]{404040}{\textbf{Laufzeit}}} & \cellcolor[HTML]{dce6f0}{\textcolor[HTML]{404040}{\textbf{Festplattenspeicher}}}\\
\midrule
\cellcolor{white}{\textcolor{black}{10.000}} & \cellcolor[HTML]{f0f8ff}{\{polars\}} & \cellcolor[HTML]{f0f8ff}{Parquet} & \cellcolor[HTML]{f0f8ff}{-} & \cellcolor[HTML]{f0f8ff}{3,3 MB} & \cellcolor[HTML]{f0f8ff}{0,4 Sek} & \cellcolor[HTML]{f0f8ff}{492,6 MB}\\
\cellcolor{white}{\textcolor{black}{10.000}} & \cellcolor[HTML]{f0f8ff}{\{data.table\}} & \cellcolor[HTML]{f0f8ff}{CSV} & \cellcolor[HTML]{f0f8ff}{Variablenauswahl} & \cellcolor[HTML]{f0f8ff}{4,4 MB} & \cellcolor[HTML]{f0f8ff}{0,7 Sek} & \cellcolor[HTML]{f0f8ff}{1,3 GB}\\
\cellcolor{white}{\textcolor{black}{10.000}} & \cellcolor[HTML]{f0f8ff}{\{arrow\}} & \cellcolor[HTML]{f0f8ff}{Parquet} & \cellcolor[HTML]{f0f8ff}{Variablenauswahl} & \cellcolor[HTML]{f0f8ff}{230,7 MB} & \cellcolor[HTML]{f0f8ff}{1,1 Sek} & \cellcolor[HTML]{f0f8ff}{492,6 MB}\\
\cellcolor{white}{\textcolor{black}{10.000}} & \cellcolor[HTML]{f0f8ff}{\{readr\}} & \cellcolor[HTML]{f0f8ff}{CSV} & \cellcolor[HTML]{f0f8ff}{Variablenauswahl} & \cellcolor[HTML]{f0f8ff}{9,3 MB} & \cellcolor[HTML]{f0f8ff}{1,2 Sek} & \cellcolor[HTML]{f0f8ff}{1,3 GB}\\
\cellcolor{white}{\textcolor{black}{10.000}} & \cellcolor[HTML]{f0f8ff}{\{arrow\}} & \cellcolor[HTML]{f0f8ff}{Parquet} & \cellcolor[HTML]{f0f8ff}{-} & \cellcolor[HTML]{f0f8ff}{229,8 MB} & \cellcolor[HTML]{f0f8ff}{1,6 Sek} & \cellcolor[HTML]{f0f8ff}{492,6 MB}\\
\cellcolor{white}{\textcolor{black}{10.000}} & \cellcolor[HTML]{f0f8ff}{\{duckdb\}} & \cellcolor[HTML]{f0f8ff}{Parquet} & \cellcolor[HTML]{f0f8ff}{-} & \cellcolor[HTML]{f0f8ff}{4,7 MB} & \cellcolor[HTML]{f0f8ff}{2,9 Sek} & \cellcolor[HTML]{f0f8ff}{492,6 MB}\\
\cellcolor{white}{\textcolor{black}{10.000}} & \cellcolor[HTML]{f0f8ff}{\{data.table\}} & \cellcolor[HTML]{f0f8ff}{CSV} & \cellcolor[HTML]{f0f8ff}{-} & \cellcolor[HTML]{f0f8ff}{3,4 GB} & \cellcolor[HTML]{f0f8ff}{3,6 Sek} & \cellcolor[HTML]{f0f8ff}{1,3 GB}\\
\cellcolor{white}{\textcolor{black}{10.000}} & \cellcolor[HTML]{f0f8ff}{\{arrow\}} & \cellcolor[HTML]{f0f8ff}{Parquet} & \cellcolor[HTML]{f0f8ff}{Datenaufteilung (Berichtsjahr)} & \cellcolor[HTML]{f0f8ff}{229,9 MB} & \cellcolor[HTML]{f0f8ff}{4,2 Sek} & \cellcolor[HTML]{f0f8ff}{486,9 MB}\\
\cellcolor{white}{\textcolor{black}{10.000}} & \cellcolor[HTML]{f0f8ff}{\{arrow\}} & \cellcolor[HTML]{f0f8ff}{Parquet} & \cellcolor[HTML]{f0f8ff}{-} & \cellcolor[HTML]{f0f8ff}{230,4 MB} & \cellcolor[HTML]{f0f8ff}{4,6 Sek} & \cellcolor[HTML]{f0f8ff}{492,6 MB}\\
\cellcolor{white}{\textcolor{black}{10.000}} & \cellcolor[HTML]{f0f8ff}{\{data.table\}} & \cellcolor[HTML]{f0f8ff}{CSV} & \cellcolor[HTML]{f0f8ff}{Datenaufteilung (Berichtsjahr)} & \cellcolor[HTML]{f0f8ff}{5,1 GB} & \cellcolor[HTML]{f0f8ff}{4,8 Sek} & \cellcolor[HTML]{f0f8ff}{1,0 GB}\\
\cellcolor{white}{\textcolor{black}{10.000}} & \cellcolor[HTML]{f0f8ff}{\{arrow\}} & \cellcolor[HTML]{f0f8ff}{Parquet} & \cellcolor[HTML]{f0f8ff}{Datenaufteilung (gleichmäßig)} & \cellcolor[HTML]{f0f8ff}{229,9 MB} & \cellcolor[HTML]{f0f8ff}{7,1 Sek} & \cellcolor[HTML]{f0f8ff}{483,5 MB}\\
\cellcolor{white}{\textcolor{black}{10.000}} & \cellcolor[HTML]{f0f8ff}{\{readr\}} & \cellcolor[HTML]{f0f8ff}{CSV} & \cellcolor[HTML]{f0f8ff}{-} & \cellcolor[HTML]{f0f8ff}{1,7 GB} & \cellcolor[HTML]{f0f8ff}{8,6 Sek} & \cellcolor[HTML]{f0f8ff}{1,3 GB}\\
\cellcolor{white}{\textcolor{black}{10.000}} & \cellcolor[HTML]{f0f8ff}{\{arrow\}} & \cellcolor[HTML]{f0f8ff}{CSV} & \cellcolor[HTML]{f0f8ff}{-} & \cellcolor[HTML]{f0f8ff}{27,4 MB} & \cellcolor[HTML]{f0f8ff}{15,7 Sek} & \cellcolor[HTML]{f0f8ff}{1,3 GB}\\
\cellcolor{white}{\textcolor{black}{10.000}} & \cellcolor[HTML]{f0f8ff}{\{data.table\}} & \cellcolor[HTML]{f0f8ff}{CSV.GZ} & \cellcolor[HTML]{f0f8ff}{-} & \cellcolor[HTML]{f0f8ff}{4,7 GB} & \cellcolor[HTML]{f0f8ff}{21,2 Sek} & \cellcolor[HTML]{f0f8ff}{421,9 MB}\\
\cellcolor{white}{\textcolor{black}{10.000}} & \cellcolor[HTML]{f0f8ff}{\{base\}} & \cellcolor[HTML]{f0f8ff}{CSV} & \cellcolor[HTML]{f0f8ff}{-} & \cellcolor[HTML]{f0f8ff}{9,7 GB} & \cellcolor[HTML]{f0f8ff}{4,9 Min} & \cellcolor[HTML]{f0f8ff}{1,3 GB}\\
\cellcolor{white}{\textcolor{black}{1.000}} & \cellcolor[HTML]{f0f8ff}{\{data.table\}} & \cellcolor[HTML]{f0f8ff}{CSV} & \cellcolor[HTML]{f0f8ff}{Variablenauswahl} & \cellcolor[HTML]{f0f8ff}{0,6 MB} & \cellcolor[HTML]{f0f8ff}{0,1 Sek} & \cellcolor[HTML]{f0f8ff}{129,8 MB}\\
\cellcolor{white}{\textcolor{black}{1.000}} & \cellcolor[HTML]{f0f8ff}{\{readr\}} & \cellcolor[HTML]{f0f8ff}{CSV} & \cellcolor[HTML]{f0f8ff}{Variablenauswahl} & \cellcolor[HTML]{f0f8ff}{2,7 MB} & \cellcolor[HTML]{f0f8ff}{0,2 Sek} & \cellcolor[HTML]{f0f8ff}{129,8 MB}\\
\cellcolor{white}{\textcolor{black}{1.000}} & \cellcolor[HTML]{f0f8ff}{\{polars\}} & \cellcolor[HTML]{f0f8ff}{Parquet} & \cellcolor[HTML]{f0f8ff}{-} & \cellcolor[HTML]{f0f8ff}{3,3 MB} & \cellcolor[HTML]{f0f8ff}{0,4 Sek} & \cellcolor[HTML]{f0f8ff}{48,1 MB}\\
\cellcolor{white}{\textcolor{black}{1.000}} & \cellcolor[HTML]{f0f8ff}{\{data.table\}} & \cellcolor[HTML]{f0f8ff}{CSV} & \cellcolor[HTML]{f0f8ff}{-} & \cellcolor[HTML]{f0f8ff}{341,8 MB} & \cellcolor[HTML]{f0f8ff}{0,9 Sek} & \cellcolor[HTML]{f0f8ff}{129,8 MB}\\
\cellcolor{white}{\textcolor{black}{1.000}} & \cellcolor[HTML]{f0f8ff}{\{duckdb\}} & \cellcolor[HTML]{f0f8ff}{Parquet} & \cellcolor[HTML]{f0f8ff}{-} & \cellcolor[HTML]{f0f8ff}{4,7 MB} & \cellcolor[HTML]{f0f8ff}{1,1 Sek} & \cellcolor[HTML]{f0f8ff}{48,1 MB}\\
\cellcolor{white}{\textcolor{black}{1.000}} & \cellcolor[HTML]{f0f8ff}{\{arrow\}} & \cellcolor[HTML]{f0f8ff}{Parquet} & \cellcolor[HTML]{f0f8ff}{Variablenauswahl} & \cellcolor[HTML]{f0f8ff}{226,2 MB} & \cellcolor[HTML]{f0f8ff}{1,2 Sek} & \cellcolor[HTML]{f0f8ff}{48,1 MB}\\
\cellcolor{white}{\textcolor{black}{1.000}} & \cellcolor[HTML]{f0f8ff}{\{arrow\}} & \cellcolor[HTML]{f0f8ff}{Parquet} & \cellcolor[HTML]{f0f8ff}{-} & \cellcolor[HTML]{f0f8ff}{225,3 MB} & \cellcolor[HTML]{f0f8ff}{1,4 Sek} & \cellcolor[HTML]{f0f8ff}{48,1 MB}\\
\cellcolor{white}{\textcolor{black}{1.000}} & \cellcolor[HTML]{f0f8ff}{\{arrow\}} & \cellcolor[HTML]{f0f8ff}{CSV} & \cellcolor[HTML]{f0f8ff}{-} & \cellcolor[HTML]{f0f8ff}{3,5 MB} & \cellcolor[HTML]{f0f8ff}{1,5 Sek} & \cellcolor[HTML]{f0f8ff}{129,8 MB}\\
\cellcolor{white}{\textcolor{black}{1.000}} & \cellcolor[HTML]{f0f8ff}{\{data.table\}} & \cellcolor[HTML]{f0f8ff}{CSV.GZ} & \cellcolor[HTML]{f0f8ff}{-} & \cellcolor[HTML]{f0f8ff}{491,7 MB} & \cellcolor[HTML]{f0f8ff}{2,7 Sek} & \cellcolor[HTML]{f0f8ff}{42,1 MB}\\
\cellcolor{white}{\textcolor{black}{1.000}} & \cellcolor[HTML]{f0f8ff}{\{data.table\}} & \cellcolor[HTML]{f0f8ff}{CSV} & \cellcolor[HTML]{f0f8ff}{Datenaufteilung (Berichtsjahr)} & \cellcolor[HTML]{f0f8ff}{510,9 MB} & \cellcolor[HTML]{f0f8ff}{2,8 Sek} & \cellcolor[HTML]{f0f8ff}{104,5 MB}\\
\cellcolor{white}{\textcolor{black}{1.000}} & \cellcolor[HTML]{f0f8ff}{\{arrow\}} & \cellcolor[HTML]{f0f8ff}{Parquet} & \cellcolor[HTML]{f0f8ff}{Datenaufteilung (Berichtsjahr)} & \cellcolor[HTML]{f0f8ff}{229,9 MB} & \cellcolor[HTML]{f0f8ff}{2,9 Sek} & \cellcolor[HTML]{f0f8ff}{53,5 MB}\\
\cellcolor{white}{\textcolor{black}{1.000}} & \cellcolor[HTML]{f0f8ff}{\{arrow\}} & \cellcolor[HTML]{f0f8ff}{Parquet} & \cellcolor[HTML]{f0f8ff}{-} & \cellcolor[HTML]{f0f8ff}{230,4 MB} & \cellcolor[HTML]{f0f8ff}{3,7 Sek} & \cellcolor[HTML]{f0f8ff}{48,1 MB}\\
\cellcolor{white}{\textcolor{black}{1.000}} & \cellcolor[HTML]{f0f8ff}{\{arrow\}} & \cellcolor[HTML]{f0f8ff}{Parquet} & \cellcolor[HTML]{f0f8ff}{Datenaufteilung (gleichmäßig)} & \cellcolor[HTML]{f0f8ff}{229,7 MB} & \cellcolor[HTML]{f0f8ff}{4,0 Sek} & \cellcolor[HTML]{f0f8ff}{50,5 MB}\\
\cellcolor{white}{\textcolor{black}{1.000}} & \cellcolor[HTML]{f0f8ff}{\{readr\}} & \cellcolor[HTML]{f0f8ff}{CSV} & \cellcolor[HTML]{f0f8ff}{-} & \cellcolor[HTML]{f0f8ff}{197,4 MB} & \cellcolor[HTML]{f0f8ff}{4,9 Sek} & \cellcolor[HTML]{f0f8ff}{129,8 MB}\\
\cellcolor{white}{\textcolor{black}{1.000}} & \cellcolor[HTML]{f0f8ff}{\{base\}} & \cellcolor[HTML]{f0f8ff}{CSV} & \cellcolor[HTML]{f0f8ff}{-} & \cellcolor[HTML]{f0f8ff}{713,4 MB} & \cellcolor[HTML]{f0f8ff}{23,3 Sek} & \cellcolor[HTML]{f0f8ff}{129,8 MB}\\
\bottomrule
\end{tabular}
\endgroup{}

}

\end{table*}%

\begin{table*}

\caption{\label{tbl-bm-reg-small}Benchmarkergebnisse für Regression (kleine Stichproben des BTP)}

\centering{

\centering\begingroup\fontsize{8}{10}\selectfont


\begin{tabular}{|>{}r||>{}l||>{}l||>{}l||>{}r||>{}r||>{}r|}
\toprule
\cellcolor[HTML]{dce6f0}{\textcolor[HTML]{404040}{\textbf{Beob.}}} & \cellcolor[HTML]{dce6f0}{\textcolor[HTML]{404040}{\textbf{Package}}} & \cellcolor[HTML]{dce6f0}{\textcolor[HTML]{404040}{\textbf{Dateiformat}}} & \cellcolor[HTML]{dce6f0}{\textcolor[HTML]{404040}{\textbf{Optimierung}}} & \cellcolor[HTML]{dce6f0}{\textcolor[HTML]{404040}{\textbf{Arbeitsspeicher}}} & \cellcolor[HTML]{dce6f0}{\textcolor[HTML]{404040}{\textbf{Laufzeit}}} & \cellcolor[HTML]{dce6f0}{\textcolor[HTML]{404040}{\textbf{Festplattenspeicher}}}\\
\midrule
\cellcolor{white}{\textcolor{black}{10.000}} & \cellcolor[HTML]{f0f8ff}{\{dplyr\}} & \cellcolor[HTML]{f0f8ff}{CSV} & \cellcolor[HTML]{f0f8ff}{Variablenauswahl} & \cellcolor[HTML]{f0f8ff}{17,6 MB} & \cellcolor[HTML]{f0f8ff}{1,3 Sek} & \cellcolor[HTML]{f0f8ff}{1,3 GB}\\
\cellcolor{white}{\textcolor{black}{10.000}} & \cellcolor[HTML]{f0f8ff}{\{arrow\}} & \cellcolor[HTML]{f0f8ff}{Parquet} & \cellcolor[HTML]{f0f8ff}{Datenaufteilung (Berichtsjahr)} & \cellcolor[HTML]{f0f8ff}{459,3 MB} & \cellcolor[HTML]{f0f8ff}{5,0 Sek} & \cellcolor[HTML]{f0f8ff}{486,9 MB}\\
\cellcolor{white}{\textcolor{black}{10.000}} & \cellcolor[HTML]{f0f8ff}{\{dplyr\}} & \cellcolor[HTML]{f0f8ff}{CSV} & \cellcolor[HTML]{f0f8ff}{-} & \cellcolor[HTML]{f0f8ff}{1,8 GB} & \cellcolor[HTML]{f0f8ff}{10,8 Sek} & \cellcolor[HTML]{f0f8ff}{1,3 GB}\\
\cellcolor{white}{\textcolor{black}{1.000}} & \cellcolor[HTML]{f0f8ff}{\{dplyr\}} & \cellcolor[HTML]{f0f8ff}{CSV} & \cellcolor[HTML]{f0f8ff}{Variablenauswahl} & \cellcolor[HTML]{f0f8ff}{4,2 MB} & \cellcolor[HTML]{f0f8ff}{0,3 Sek} & \cellcolor[HTML]{f0f8ff}{129,8 MB}\\
\cellcolor{white}{\textcolor{black}{1.000}} & \cellcolor[HTML]{f0f8ff}{\{dplyr\}} & \cellcolor[HTML]{f0f8ff}{CSV} & \cellcolor[HTML]{f0f8ff}{-} & \cellcolor[HTML]{f0f8ff}{206,9 MB} & \cellcolor[HTML]{f0f8ff}{6,4 Sek} & \cellcolor[HTML]{f0f8ff}{129,8 MB}\\
\cellcolor{white}{\textcolor{black}{1.000}} & \cellcolor[HTML]{f0f8ff}{\{arrow\}} & \cellcolor[HTML]{f0f8ff}{Parquet} & \cellcolor[HTML]{f0f8ff}{Datenaufteilung (Berichtsjahr)} & \cellcolor[HTML]{f0f8ff}{452,5 MB} & \cellcolor[HTML]{f0f8ff}{6,8 Sek} & \cellcolor[HTML]{f0f8ff}{53,5 MB}\\
\bottomrule
\end{tabular}
\endgroup{}

}

\end{table*}%


\printbibliography



\end{document}
